\documentclass[ecta,nameyear,draft]{econsocart}
\usepackage{amsmath,amsfonts,amssymb,amsthm,booktabs,array,longtable}
\usepackage[T1]{fontenc}
\usepackage{mathpazo,microtype,needspace,xurl}
\RequirePackage[colorlinks,citecolor=blue,linkcolor=blue,urlcolor=blue]{hyperref}
\startlocaldefs
\newtheorem{theorem}{Theorem}
\newtheorem{proposition}[theorem]{Proposition}
\newtheorem{lemma}[theorem]{Lemma}
\newtheorem{corollary}[theorem]{Corollary}
\newcommand{\E}{\mathbb E}
\newcommand{\R}{\mathbb R}
\endlocaldefs
\setlength{\emergencystretch}{3em}
\makeatletter\def\copyright@text{Prepared for referee review}\makeatother

\begin{document}
\begin{frontmatter}
\title{Neural Bellman Operators}
\runtitle{Neural Bellman Operators}
\begin{aug}
\author[id=au1,addressref={add1}]{\fnms{Qian}~\snm{QI}\ead[label=e1]{qiqian@pku.edu.cn}}
\address[id=add1]{Peking University}
\end{aug}
\begin{abstract}
This paper develops Neural Bellman Operators for policy evaluation and feasible improvement in controlled economies. Centered continuation errors connect economic transport to full-policy loss. A deterministic own-future ReLU construction retains the feasible actions attaining its labels and supplies a one-sided Bellman certificate. We extend this construction to policies implemented from acquired states: numerical selection, state uncertainty, endogenous capacity and action quantization enter one feasible-execution budget without assuming actor continuity. An oracle-explicit coupled economy makes the covering, continuous-law integration and repair operations executable. Direct expected-cost comparisons evaluate the actual witness policies under a prespecified simultaneous uncertainty account, and a catalogue-adoption theorem converts these comparisons into economic choices. Complete work-to-bound curves distinguish sharper certification from first-target work. The evidence identifies tolerance-based policy-cost proximity and implemented constrained guarantees while preserving unfavorable conventional comparisons. Controlled-diffusion, recursive-preference and game formulations retain their application-specific conditions.
\end{abstract}
\begin{keyword}
\kwd{Dynamic programming}\kwd{Neural approximation}\kwd{Policy certification}
\kwd{Recursive utility}\kwd{Computational economics}
\end{keyword}
\end{frontmatter}

\section{Introduction}\label{sec:introduction}
An economic decision uses a continuation value because its consequences extend beyond the current period. A learned continuation is useful only to the extent that the decisions made with it are accurate. Prediction loss, a Bellman residual at sampled states, and the value of the policy actually implemented are different objects. Their distinction is particularly important when a counterfactual changes future policies or when a recursive preference evaluates the continuation generated by those policies rather than the optimum.

This paper develops Neural Bellman Operators (NBOs) as a policy-evaluation and feasible-improvement method. An operator takes economic primitives, a current policy, and computational resources as inputs. It returns a continuation, a feasible policy, and an economic error account, or it records that the requested account was not obtained within its cap. The continuation is evaluated against its own fitted future. Current actions, policy-specific future values, numerical execution, and the cost of verification enter the same account. The construction is not a proposal to rename an arbitrary neural predictor as a Bellman operator.

There are four parts to the argument. First, only action-dependent continuation errors affect a current decision. Their oscillation, rather than their absolute level, controls loss. A transport bound identifies when an anchor continuation remains valid after the future economy changes. These results connect reuse and refresh decisions to economic tolerances. They also preserve the original formulation in controlled diffusions, recursive preferences, endogenous preferences, and games, subject to the distinct hypotheses of those applications.

Second, a signed residual certificate controls a complete returned policy. Its upper bound contains the widths of the residual enclosures, the suboptimality of the selected actions, and the terminal enclosure. It is invariant to additive date-specific shifts of the fitted continuation. This leads to a useful separation: a policy can be accurate even though a coarse verifier does not establish its accuracy. We prove that residual verification can be refined while retaining the original network, the original deployed actor, and its original allowance. The resulting bound measures improved knowledge of the same policy, not a policy improvement disguised as additional verification.

Third, certification of each policy relative to the optimum does not rank two policies. We construct their cost difference directly. Under expected-cost preferences, an own-continuation Bellman telescoping identity produces an unbiased policy-cost score with a support width governed by the signed certificate. Coupling innovations across policies and enclosing every simulated path gives simultaneous finite-sample intervals for the difference in the original continuous-innovation economy. The statistical assertion is conditional on frozen policies and the stated independent simulation model; outward arithmetic and unresolved action-cell boundaries have their own explicit treatment.

Fourth, a deterministic neural construction connects primitive moduli to a complete policy tolerance. A minimum of Lipschitz cones, realized exactly by ReLU gates, fits each date's own-future Bellman labels without inflating its Lipschitz bound when those labels are inexact. The state, action, and arithmetic budgets are specified before construction. This yields finite termination for a defined neural backend rather than a claim about arbitrary nonlinear training. A new full matched catalogue executes the same resolution ladder for this backend and a conventional spline, including every unsuccessful rung and three isolated timing repetitions.

The action that attains a Bellman label supplies additional information. We retain its index through the ReLU minimum circuit and deploy that action when its cone is selected. A one-sided policy sandwich then avoids charging the separate actor and critic accounts twice. Both the nearest-node neural ablation and the spline receive a strengthened certificate. Their complete matched frontiers isolate certificate sharpness from first-target work and actual policy cost.

A feasible policy must also survive the information and arithmetic used to implement it. We give a uniform acquisition-aware theorem that permits discontinuous witness selection, uncertain measured states and capacity-dependent repair. A coupled two-state investment economy realizes all covering, integration and repair operations explicitly, with a conventional fitted-value comparator. The original scalar witness policies receive direct expected-cost comparisons at every common first crossing, under a tolerance and simultaneous error account fixed before execution. This establishes proximity within the declared decision band, not a signed neural ranking. A catalogue-choice theorem makes the result usable when implementation charges differ. The complete resource curves and an identical native-envelope control isolate what the construction, the certificate and the neural representation each contribute.

The numerical implementation starts from economic primitives. In the scalar investment economy, it fits a fresh ReLU continuation backward against targets using its own previously fitted future network. In the two-capital economy, it trains hidden locations in a convex squared-ReLU continuation, with continuous simplex actions and continuous uniform innovations. The second experiment includes a flexible nonnegative ridge dictionary as well as a quadratic projection. Each comparator constructs its own future targets and uses the same final economic criterion. Neither conventional values nor inherited neural weights are used as labels for the neural services.

We report the complete 150-service catalogue, not only successful candidates. In the original capped scalar services the neural procedure reaches the all-state target $0.06$ in all twelve cases and $0.04$ in six, whereas the spline reaches both targets in every case. In the coupled services the neural procedure reaches $0.25$ in six cases and $0.10$ in five; the flexible ridge reaches both in all six. These statements describe the original service, including its original verification cap.

A subsequent diagnostic holds both the neural weights and the deployed policies fixed. Refining only residual verification certifies all six previously unresolved scalar policies at $0.04$ and the previously unresolved coupled policy at $0.10$. The original failures are not relabeled. The diagnostic shows that they do not, by themselves, identify a failure to construct an accurate neural policy. It does not identify the population success probability of the optimizer, establish a new end-to-end runtime, or change any policy's cost.

The comparative findings are equally integral to the paper. The original scalar signed policy comparisons frequently favor the spline. The new coupled direct intervals all contain zero and therefore do not order the policies. The precision controller genuinely escalates after rejected low-precision proposals, but its complete recorded work does not improve on fixed binary64 in this experiment. These findings distinguish the certified capability of the construction from a comparative advantage that would require additional evidence.

The paper retains its original economic subject and full theory. The exposition now follows the economic problem, the operator, the policy theorem, and the evidence. The complete previous article, proof supplement, and economic applications are materialized unchanged within the current source tree as preserved companions; the preservation map records their exact identities. No historical theorem or application is deleted to obtain a shorter active exposition. The present article and new technical supplement state and prove the additional results used by the new evidence.

\subsection{Relation to the literature}
Neural Bellman Operators sit within approximate dynamic programming, fitted-value methods, and policy iteration, not outside them. Regression-based continuation methods and simulation-based dynamic programming supply important conventional benchmarks \citep{judd1998,powell2011,longstaff2001,ernst2005}. Actor--critic and approximate policy-iteration methods separate evaluation and improvement \citep{konda1999,scherrer2015}; safe-improvement arguments likewise distinguish a proposed action from a verified gain \citep{kakade2002,pirotta2013}. We use these distinctions to organize the economic object returned by the algorithm.

Neural methods for differential equations provide ways to approximate values and their derivatives \citep{sirignano2018,han2018,hornik1990}. Such approximation does not on its own verify a returned dynamic policy or its numerical execution. Our certificates can be applied to neural or conventional continuations. The neural-specific construction lies in fresh own-future training and, in the coupled experiment, learned hidden locations. Its comparative value must be isolated against flexible conventional methods, not inferred from the certificate's validity.

The general economic formulation admits recursive utility \citep{epsteinzin1989,duffieepstein1992}. The compact-domain policy theorem uses monotonicity and cash invariance, while the original Gaussian entropic construction has separate moment and spectral conditions. These are different bridges, not interchangeable assumptions. The direct paired-cost identity below applies to conditional expectation; it is not asserted for a nonlinear certainty equivalent. The confidence step uses the empirical Bernstein inequality of \citet{maurer2009}. Our additional ingredients are the policy-specific telescoping score, its certificate-derived support, and the numerical enclosure of the original continuous-law comparison.

\section{The controlled economy and the operator}\label{sec:economy}
\subsection{Policy-specific continuation}
The original continuous-time economy has an admissible control $a$ and state
\begin{equation}\label{eq:diffusion}
 dS_t=b(t,S_t,a_t)\,dt+\sigma(t,S_t,a_t)\,dW_t.
\end{equation}
The feasible correspondence, invariant or stopped domain, and exit and terminal rewards are economic primitives. For recursive evaluation, the continuation of a fixed policy is represented, on its specified horizon, by
\[
 Y_s^a=g(\tau,S_\tau)+\int_s^\tau f(v,S_v,a_v,Y_v^a)\,dv
                 -\int_s^\tau Z_v^a\,dW_v.
\]
The value is $V^a(t,s)=Y_t^a$ under the well-posedness and comparison hypotheses in the retained controlled-economy analysis. In particular, the recursive argument is the policy's own continuation. Replacing it by the optimal continuation would change the policy-evaluation problem.

The applications companion retains the full boundary specifications, infinite-horizon transversality conditions, utility-domain restrictions, temporal-self equations, and game conditions. A finite-domain residual calculation does not establish any of these by itself. In a game, certifying a specified deviation problem is not the same as establishing the equilibrium fixed point. In continuous time, the controlled generator contains state derivatives; freezing a critic during an actor update does not suppress legitimate derivatives of the represented state-dependent function.

For the constructive implementation and the main policy theorem, consider dates $t=0,\ldots,T$, compact state space $X$, and feasible correspondence $A_t(x)$. States evolve as
\begin{equation}\label{eq:transition}
 X_{t+1}=F_t(X_t,a_t,Z_{t+1}).
\end{equation}
Let $c_t$ be cost, $g$ terminal cost, $\beta_t>0$, $B_0=1$, and $B_t=\prod_{j<t}\beta_j$. We use a cost-minimization convention in the discrete implementation. It is equivalent to changing signs in an additive reward formulation, but does not replace the original recursive-utility assumptions by a sign convention.

Write
\begin{align}\label{eq:bellman}
 Q_t(v;x,a)&=c_t(x,a)+\beta_t\rho_t\bigl(v(F_t(x,a,Z_{t+1}))\bigr),\\
 \mathcal T_tv(x)&=\inf_{a\in A_t(x)}Q_t(v;x,a),\qquad V_T^*=g.\notag
\end{align}
The conditional evaluations $\rho_t$ are monotone and cash invariant on a domain containing every value and comparison function used below. Feasible measurable selectors exist; compactness, continuity of primitives, and continuity of the evaluated integrals are sufficient conditions in the experiments. The Bellman recursion is valid against all admissible adapted policies. For expectation, $J^\pi$ is expected discounted cost. For a recursive evaluation it denotes the value generated by that evaluation's own recursion.

\subsection{An executable NBO service}
An implemented operator maps economic data and a current policy to
\[
 (\widehat f,\widehat\pi,\mathcal C,W)\quad\hbox{or}\quad
 (\mathrm{unresolved},W,\mathrm{diagnostic}),
\]
where $\mathcal C$ is an economic certificate and $W$ records the work used. Its stages are target construction, fitting, feasible action selection, independent certification, and durable storage. Critic fitting freezes its future inputs; action selection freezes the resulting critic. In the fresh backward construction, date $t$ uses the already fitted $\widehat f_{t+1}$, not an externally solved optimal continuation. The terminal target is the primitive $g$.

A finite optimizer schedule is an algorithm, not a convergence theorem. A service that has not met its target at its cap records failure, with the complete cost of its attempts. The retained factor-training analysis provides a finite constructive theorem for a separate structured class. The nonlinear schedules below instead have an explicit, auditable capped-service estimand.

For a fixed catalogue $\mathcal K$ and target $\varepsilon$, define
\begin{equation}\label{eq:estimand}
 H_m(k,\varepsilon)=\mathbf1\{\text{method $m$ stores a valid certificate
 at $\varepsilon$ within its cap on object $k$}\}.
\end{equation}
The empirical object is the full array of $H_m$, costs to the first qualifying checkpoint, and capped costs for failures. Its arithmetic average is a catalogue fraction. A probability over random initializations would require a specified law and a sampling design for that different estimand. A post-review refinement of a failed object's verifier is a new diagnostic and does not change \eqref{eq:estimand} for the original service.

\section{Continuation accuracy, transport, and policy loss}\label{sec:theory}
\subsection{Centered errors and economic reuse}
Consider a current objective $r_\theta(s,a)+S_\theta(T_\theta(s,a))$ to be maximized. Let $e=\widehat S_0-S_0$ be anchor approximation error and $\Delta_\theta=S_\theta-S_0$ the change in the actual future continuation. Suppose the implemented action is $\delta$-optimal for $r_\theta+\widehat S_0\circ T_\theta$.

\begin{proposition}[Centered decision and transport bounds]\label{prop:center}
The current decision loss satisfies
\begin{equation}\label{eq:osc}
 L_\theta(s,\widehat a)\leq\delta+
 \operatorname{osc}_{A_\theta(s)}(e\circ T_\theta)+
 \operatorname{osc}_{A_\theta(s)}(\Delta_\theta\circ T_\theta).
\end{equation}
For a common finite action menu with positive weights $p_j$, set
$R_p=\sum_jp_j(e_j-\sum_ip_ie_i)^2$ and
$C_p=\max_{i\ne j}(p_i^{-1}+p_j^{-1})$. If $|\Delta_\theta|\leq b_\theta$ on the queried postdecision set, then
$L_\theta\leq\delta+\sqrt{C_pR_p}+2b_\theta$.

If common-innovation future transitions have Lipschitz constants $L_k$ and perturbation bounds $\epsilon_{\theta,k}$, and future rewards have Lipschitz constants $H_k$ and perturbations $\eta_{\theta,k}$ on an invariant domain, put $d_0=0$, $d_{k+1}=L_kd_k+\epsilon_{\theta,k}$. For nonnegative reward weights $w_k$,
\begin{equation}\label{eq:transport}
 |S_\theta-S_0|\leq b_\theta:=\sum_kw_k(\eta_{\theta,k}+H_kd_k).
\end{equation}
\end{proposition}

The proof compares the true and fitted objective at the two selected actions and couples the future states. The weighted finite-menu bound is Cauchy--Schwarz applied to $e_i-e_j$. Details are in the technical supplement. An additive continuation error common to all actions is irrelevant to the decision. A finite menu, however, does not certify a continuous action set merely because it includes both observed actions.

For recursive maps with sup-norm Lipschitz constants $\beta_k$ on a common invariant utility domain, a one-step perturbation $\nu_k$ propagates as $b_k\leq\nu_k+\beta_kb_{k+1}$. This is the recursive counterpart of \eqref{eq:transport}. Its use requires the common domain; it cannot cross a singular utility boundary by extending a local constant. These transport bounds retain the original refresh logic: use an anchor while its economic allowance covers approximation and future change, and recertify or recompute when that condition is exhausted.

\subsection{A direct all-state policy theorem}
For continuous fitted functions $f_t$ and a feasible measurable policy $\pi$, suppose the verifier supplies, uniformly in $x$,
\begin{align}\label{eq:residual}
 \ell_t\leq f_t(x)-\mathcal T_tf_{t+1}(x)&\leq u_t,\qquad t<T,\\
 Q_t(f_{t+1};x,\pi_t(x))-\mathcal T_tf_{t+1}(x)&\leq\eta_t,\notag\\
 \ell_T\leq f_T(x)-g(x)&\leq u_T.\notag
\end{align}
These are enclosures for the actual coefficients and actual returned actions. Neither an optimizer's success flag nor a residual on sampled states replaces them.

\begin{theorem}[Centered full-policy certificate]\label{thm:policy}
Under the conditions of Section~\ref{sec:economy} and \eqref{eq:residual}, every initial state satisfies
\begin{equation}\label{eq:gap}
 0\leq J^\pi_0(x)-V_0^*(x)\leq
 G(f,\pi):=\sum_{t<T}B_t(u_t-\ell_t+\eta_t)+B_T(u_T-\ell_T).
\end{equation}
More specifically,
\begin{align}\label{eq:brackets}
 V_0^*(x)&\geq f_0(x)-\sum_{t<T}B_tu_t-B_Tu_T,\\
 J_0^\pi(x)&\leq f_0(x)+\sum_{t<T}B_t(\eta_t-\ell_t)-B_T\ell_T.\notag
\end{align}
Replacing $f_t$ by $f_t+d_t$, with date-specific constants $d_t$, leaves the bound unchanged after the residual endpoints are shifted by $d_t-\beta_td_{t+1}$ and the terminal endpoints by $d_T$.
\end{theorem}

\begin{proof}
Monotonicity and cash invariance imply
$\mathcal T_t(v+c)=\mathcal T_tv+\beta_tc$ and the analogous identity for the fixed-policy map. Starting from $g\geq f_T-u_T$, backward induction gives the first inequality in \eqref{eq:brackets}. Starting from $g\leq f_T-\ell_T$ and using
$Q_t(f_{t+1};x,\pi_t(x))\leq f_t(x)-\ell_t+\eta_t$
gives the second. Subtraction proves the upper bound; optimality proves the lower bound. The shifted residuals have unchanged widths and unchanged action gaps.
\end{proof}

The theorem certifies a full adapted policy, not a finite action menu or a grid of initial states. With expected cost its proof is a Bellman telescoping argument; for a monotone cash-invariant recursive evaluation it is an order argument. The numerical experiment is not needed for its validity. Its numerical usefulness depends on producing small, valid enclosures.

\begin{corollary}[Retained-policy recertification]\label{cor:retain}
Fix the functions $f_t$, the deployed policy $\pi$, valid actor allowances $\eta_t^{\rm old}$, and a terminal enclosure. Let a second verifier return valid residual enclosures $[\ell_t^{\rm new},u_t^{\rm new}]$ for those same functions. Then
\begin{equation}\label{eq:retained}
 J_0^\pi-V_0^*\leq
 \sum_{t<T}B_t(u_t^{\rm new}-\ell_t^{\rm new}+\eta_t^{\rm old})
          +B_T(u_T^{\rm old}-\ell_T^{\rm old}).
\end{equation}
No retraining or change in deployed actions is required. If desired, intersecting two valid residual enclosures makes the new bound no larger than the old one.
\end{corollary}

The original actor allowance is valid on the entire state domain, so a finer critic verification mesh does not invalidate it. Conversely, replacing it by an allowance for a newly optimized fine-grid actor would certify a different policy. The implementation records the original actor hash and explicitly discards any fine-grid action proposals generated internally by the verifier. Those proposals' work and transient storage are still charged.

The corollary does not promise that arbitrary precision is attainable by mesh refinement. Nodal residual variation, the retained actor allowance, and the terminal enclosure can leave a positive floor. Nor is the corollary a claim that additional verification improves welfare. It distinguishes the accuracy of a fixed economic decision rule from the sharpness of our knowledge of that accuracy.

\subsection{Construction beyond post-training verification}
The original trainable-factor construction remains a constructive part of NBO. For a quadratic own-policy continuation $x'P_tx+c_t$, a square-activation network $\|W_tx\|^2/d+\widehat c_t$ trains every entry of $W_t$. Its population derivative loss and gradient are
\[
 \mathcal L_t(W;P)=\tfrac14\|W'W-dP\|_F^2,
 \qquad \nabla_W\mathcal L_t=W(W'W-dP).
\]
If $\sigma_{\min}(W)\geq s>0$, then
$\|W'W/d-P\|_F\leq\|\nabla_W\mathcal L_t\|_F/(ds)$.
Together with the full-vector actor defect, this gives a policy-error account for a jointly trained hidden factor. The own-policy coefficient is re-evaluated for the returned policy, not silently borrowed from an earlier training policy.

The primitive-budget analysis in the retained article chooses a feasibility witness, coefficient envelopes, rank margins, and update budgets before training. In the Gaussian entropic extension it also verifies positive exponential-moment margins. Under its stated spectral hypotheses, a scaled orthogonal initialization and bounded gradient steps give the finite training construction proved there. The finite matrix and Gaussian formulas are not a nonlinear dimension-scaling experiment. They also do not imply that Adam or L-BFGS has a deterministic residual cap for a general squared-ReLU network. We preserve both results and keep the two logical statements separate.

\section{From economic primitives to nonlinear certificates}\label{sec:algorithm}
\subsection{Backward fitting and selected policies}
For each method, initialize a fresh terminal fit from $g$. At date $t$, construct targets from $\mathcal T_tf_{t+1}$ using that method's own already fitted future. Fit the current continuation under the frozen schedule. Then propose feasible actions, enclose the Bellman infimum and the selected-action gap, and validate residual coverage. Store a checkpoint before reporting its first target crossing. Continue through the declared ladder only as specified by the original study. The record includes all attempted resolutions, unsuccessful target queries, target-generation work, and storage operations.

The scalar network has 63 ReLU features. Its training uses 513 states, a 257-point target action construction, 400 Adam steps at learning rate $5\times10^{-5}$, and a final output fit. The scalar verifier uses its own finer action and state ladders. The conventional spline is reconstructed from the same primitives with its own ladder and targets. A conventional target solver is not used to label a neural continuation.

The coupled neural continuation is
\begin{equation}\label{eq:network}
 f_t(x)=a_t+b_t'x+\sum_{j=1}^{48}c_{tj}(w_{tj}'x+v_{tj})_+^2,
 \qquad c_{tj}\geq0.
\end{equation}
Hidden locations and weights are trained. Nonnegative coefficient initialization and a finite L-BFGS schedule of 300 iterations act on fresh own-future targets. The ridge comparator instead fixes a 96-feature squared-hinge dictionary and fits nonnegative coefficients to its own targets; zero coefficients may be pruned. Its dictionary seed is fixed by date, so repeated training-seed labels do not create independent random ridge representations. A quadratic projection is retained as a restricted comparator, not treated as the strongest conventional alternative.

\subsection{All-state and continuous-action verification}
At each state node, outward arithmetic encloses the actual network value and the action objective. In the coupled economy, convexity in actions and a positive strong-convexity constant yield a continuous-simplex primal--dual gap. A projected optimizer merely proposes the action; its stopping flag is not the certificate. State interpolation then extends the signed residual bounds off-grid. The nearest-node deployed actor receives a separate, uniform allowance for using its node action at the actual state.

Uniform innovations are integrated through positive-part moments, rather than replaced by a sampled innovation law. The elementary identity for a uniform variable $U$ on $[-h,h]$ is
\[
 \E(y+U)_+^2=\frac{(y+h)_+^3-(y-h)_+^3}{6h},\qquad h>0.
\]
Successive integration gives the corresponding formula for independent coordinates and affine neuron arguments; zero coefficients are handled by the lower-dimensional formula. Interval evaluation encloses cancellation and rounding. The supplement specifies the mesh and action bounds used in the inherited verifier.

The residual account distinguishes a nodal enclosure from state coverage, actor error, and terminal error. A nodal enclosure contains the effects of representation, fitting, Bellman optimization, and arithmetic; it is not an identified decomposition into a pure approximation error and a pure optimizer error. In particular, subtracting a training loss from the residual does not estimate either component. The diagnostic reports every separately available allowance without counting nested subcomponents twice.

\subsection{Work, dimension, and execution}
For tensor state, action, and innovation covers of sizes $K_x,K_a,m$, a direct verifier pays work of order
\begin{equation}\label{eq:work}
 \sum_{t<T}K_xK_a\{m(C_F+P_{t+1})+C_\rho\}
       +\sum_{t\leq T}K_xP_t+W_{\rm storage},
\end{equation}
where $P_t$ counts network evaluation work. Actor storage is of order $TK_xd_a$. Analytic integration replaces $m$ by the work of the moment formula, not by zero. With tensor covers, $K_x$ scales as $h_x^{-d_x}$; analogous factors arise for action and innovation discretizations. A dimension-explicit certificate is therefore not a dimension-free algorithm.

If a Lipschitz innovation approximation contributes at most $\zeta_t$ to $Q_t$, widening residual endpoints by $\zeta_t$ and the actor allowance by $2\zeta_t$ adds $4\sum_tB_t\zeta_t$ to \eqref{eq:gap}. This is a discretization account, not permission to replace a continuous law without charging its error. Independent policy queries may have their own scenario-tree or simulation cost and are recorded separately.

The complete service includes own-future targets, fitting, unsuccessful refinement, action extraction, verification, applicable policy queries, serialization, and durable checkpoint storage. Initialization and outer-process time are retained separately. New local diagnostic clocks are not added to old GitHub-hosted clocks to manufacture an end-to-end estimate. All reported arithmetic certificates concern stored binary coefficients with outward enclosures. The inherited execution account treats rounded state acquisition and transition/cost evaluation separately from ideal policy loss.

\section{Direct comparisons of economic policy costs}\label{sec:paired}
A smaller upper bound on $J^\pi-V^*$ need not mean a smaller $J^\pi$. We therefore compare the neural and ridge policies directly at common states, using common innovations and the expectation version of \eqref{eq:bellman}.

For method $m$, let $X^m$ follow its own fixed policy and define
\begin{equation}\label{eq:score}
 Z_m=f_0^m(x_0)+\sum_{t<T}B_t
  \{Q_t(f_{t+1}^m;X_t^m,\pi_t^m(X_t^m))-f_t^m(X_t^m)\}
       +B_T\{g(X_T^m)-f_T^m(X_T^m)\}.
\end{equation}
This is a control-variate score for actual policy cost, not the difference of two regret bounds.

\begin{theorem}[Paired economic score and support]\label{thm:paired}
For fixed policies and integrable continuations under conditional expectation,
$\E Z_m=J^{\pi^m}_0(x_0)$. If method $m$ satisfies \eqref{eq:residual}, then every trajectory satisfies
\begin{align}\label{eq:support}
 L_m&\leq Z_m\leq U_m,\\
 L_m&=f_0^m(x_0)-\sum_{t<T}B_tu_t^m-B_Tu_T^m,\notag\\
 U_m&=f_0^m(x_0)+\sum_{t<T}B_t(\eta_t^m-\ell_t^m)-B_T\ell_T^m.\notag
\end{align}
Consequently $Z_N-Z_R$ has expectation $J^{\pi^N}-J^{\pi^R}$ and support
$[L_N-U_R,U_N-L_R]$, whose width is $G_N+G_R$. These statements hold under any coupling of the two methods that preserves each policy's innovation law.
\end{theorem}

\begin{proof}
Conditional expectation gives $\E Q_t(f_{t+1}^m;X_t^m,\pi_t^m)=\E c_t+\beta_t\E f_{t+1}^m(X_{t+1}^m)$. Summing telescopes all continuation terms. The selected-action gap is between zero and $\eta_t^m$, so each centered summand in \eqref{eq:score} is between $-u_t^m$ and $\eta_t^m-\ell_t^m$. The terminal summand is between $-u_T^m$ and $-\ell_T^m$. Summing proves the support statement and its width.
\end{proof}

To evaluate the original continuous law, simulate uniform innovation bins and enclose all within-bin innovations. Interval propagation includes every actor that can be selected when a state enclosure crosses a nearest-node boundary. Suppose the resulting endpoint pair satisfies $A_i\leq Z_{N,i}-Z_{R,i}\leq D_i$ and lies in a known support interval of width $B$. For independent draws of bins, the lower and upper endpoints are bounded independent observations. Let $\overline A,\overline D$ and $s_A^2,s_D^2$ be their sample means and variances. For $q$ predeclared comparisons, set
\begin{equation}\label{eq:bernstein}
 r(s^2)=\sqrt{\frac{2s^2\log(4q/\alpha)}{n}}
           +\frac{7B\log(4q/\alpha)}{3(n-1)}.
\end{equation}
Applying the one-sided empirical Bernstein inequality to the lower and upper endpoints separately, then taking a union bound, gives simultaneous coverage at least $1-\alpha$ for
\begin{equation}\label{eq:ci}
 [\overline A-r(s_A^2),\ \overline D+r(s_D^2)]
\end{equation}
over all $q$ comparisons. Mean, variance, logarithm, and square-root calculations are outward enclosed in the implementation; intersecting with the deterministic support preserves coverage. The proof and numerical details are in the supplement.

The coverage statement is conditional on frozen training outputs and an independent simulation model. Fixed PCG64 streams make this implementation reproducible; a fixed stream is not itself a proof of probabilistic independence. No policy is selected using these simulations. Intervals crossing zero are reported as unresolved, not as equivalence or as evidence of a superior neural policy.

\section{Investment economies and the complete evidence}\label{sec:evidence}
\subsection{Design and provenance}
The original science design has twelve scalar neural services, twelve scalar spline services, six coupled neural services, six coupled ridge services, six coupled quadratic services, and 36 services for each of three precision rules. The successful source-bound execution contains all 150 records. The original first run, with four nonfinite-network exceptions, is retained. Its numerical inverse-softplus initialization was subsequently replaced by an algebraically equivalent stable expression; economic primitives, seeds, targets, budgets, and comparators were unchanged. We do not treat the failed first run as nonexistent or combine its clocks with the successful run.

The R44 protocol was committed before the additional diagnostics were executed. It is a post-review protocol, not an original preregistration. It selects the seven unresolved tight-target neural objects by their already observed outcomes and freezes their weights and policies. The paired comparison uses all six coupled pairs and all four originally declared states. A source manifest binds the original inputs, the diagnostic code, outcomes, and publication sources. Every original record and all five original science-source hashes were replayed.

\subsection{Scalar and coupled primitives}
The scalar experiment retains the original bounded investment state, price and risk cells, three declared training seeds, continuous feasible action interval, and the independent own-policy evaluator. Prices are one and four. Both methods construct their own continuations from the same primitive problem. The retained evidence and code specify every scalar primitive; the full economic derivation remains in the original investment section and applications.

In the coupled economy, $x\in[0,1]^2$, $T=4$, $\beta=15/16$, and
$A=\{a\geq0:a_1+a_2\leq1/8\}$. With independent $z_j\sim U[-1/32,1/32]$,
\[
 F_j(x,a,z)=\tfrac18+\tfrac12x_j+\tfrac18x_{3-j}
           +\tfrac1{16}x_j(1-x_{3-j})+a_j+z_j.
\]
This transition maps the stated state-action-innovation domain into the state domain. Define
\[
 s(x)=\sum_{j=1}^2(x_j-11/16)^2
       +2(1/2-x_1-x_2)_+^2+\tfrac14(x_1-x_2)^4.
\]
The stage cost is $s(x)+p\sum_ja_j^2+\tfrac12a_1a_2+4\sum_ja_j^4$; terminal cost doubles the squared-target term in $s$ and retains its remaining terms. Prices $p\in\{1,4\}$ penalize costly investment. The shared simplex represents a common resource constraint, the transition couples the two capital stocks, and the state cost penalizes aggregate shortage and imbalance. This is a stylized controlled economy, not an empirically calibrated welfare exercise.

For both coupled methods the final all-state verification ladder is $N=16,32,64,128$. The action objective has strong-convexity lower bound $2p-1/2>0$. Uniform positive-part moments integrate continuous uncertainty analytically. The finest original actor has $4(129)^2\cdot2=133{,}128$ stored scalar actions. This actor cost is distinct from continuation-network storage.

\subsection{Original attainment and comparative outcomes}
\begin{table}[htbp]
\centering\small
\caption{Original capped-service target attainment}
\label{tab:attainment}
\begin{tabular}{lrr}
\toprule
Method & All-state target & Certified\\
\midrule
Scalar neural & 0.06 & 12/12\\
Scalar neural & 0.04 & 6/12\\
Scalar spline & 0.06 & 12/12\\
Scalar spline & 0.04 & 12/12\\
Coupled neural & 0.25 & 6/6\\
Coupled neural & 0.10 & 5/6\\
Coupled ridge & 0.25 & 6/6\\
Coupled ridge & 0.10 & 6/6\\
Coupled quadratic & 0.25 & 0/6\\
Coupled quadratic & 0.10 & 0/6\\
\bottomrule
\end{tabular}
\end{table}

Table~\ref{tab:attainment} reports the original capped-service outcomes. The original scalar neural bound is larger than its matched spline bound in every one of the twelve pairs. At the first common $0.06$ crossing, 29 of 48 signed neural-minus-spline intervals have positive lower endpoints, none has a negative upper endpoint, and 19 contain zero. At $0.04$, the six common pairs give 17 positive intervals, none negative, and seven overlaps. All lie within the prespecified descriptive margin $\pm0.005$. This establishes numerical closeness at those states but is not an equivalence design; where signs are identified, they favor the spline.

The coupled ridge reaches both targets in all six services. It has the smaller final gap in four pairs and the neural method in two. At $0.25$ the ridge reaches the target faster in all six original observations. At $0.10$ it is faster in four of the five pairs in which both pass, while the neural method is faster in one and fails in the remaining pair. These singleton nonlinear timing observations are descriptive. The ridge continuation has storage comparable to the neural continuation, so its strength cannot be dismissed as a comparison with an artificially restricted quadratic.

\subsection{Diagnosing tight-target failures without changing policies}
\begin{table}[htbp]
\centering\small
\caption{Additional verification of unchanged neural policies}
\label{tab:recertification}
\begin{tabular}{lrrrr}
\toprule
Service & Original bound & Refined bound & Target & Attained\\
\midrule
041 & 0.04092128 & 0.03068472 & 0.04 & Yes\\
089 & 0.04118465 & 0.03090917 & 0.04 & Yes\\
120 & 0.04112077 & 0.03085353 & 0.04 & Yes\\
124 & 0.04093160 & 0.03070697 & 0.04 & Yes\\
128 & 0.15223570 & 0.09775755 & 0.10 & Yes\\
145 & 0.04136567 & 0.03108538 & 0.04 & Yes\\
147 & 0.04101560 & 0.03077643 & 0.04 & Yes\\
\bottomrule
\end{tabular}
\end{table}

For the six unresolved scalar policies, only the residual state grid is refined from $N=2048$ to $N=4096$, retaining the original action verification resolution $A=1024$. For the coupled policy, residual verification uses $N=256$ instead of $128$. The original actor allowances and terminal enclosures remain in \eqref{eq:retained}; all original action arrays and networks are hash-checked unchanged. Fine-grid actor proposals generated internally are discarded, not deployed.

Table~\ref{tab:recertification} gives all seven results. The retained scalar policy bounds become approximately $0.03068$--$0.03109$, below $0.04$. The retained coupled policy bound becomes $0.09775756$, below $0.10$. There are zero training updates and zero newly stored deployed-action scalars. The coupled diagnostic still pays for a fine verification mesh with 66,049 states, including transient fine-grid proposals; retaining the original actor does not make verification free.

The original coupled decomposition is informative. Its discounted state-cover contribution to nonterminal residual widths is about $0.07268224$, its actor allowance is $0.03287249$, its remaining nodal residual enclosure is $0.04574738$, and its terminal contribution is $0.00093360$. These sum to the original bound up to displayed rounding. The extra mesh quarters the principal second-order coverage term while retaining the actor and terminal terms. The resulting crossing does not establish that representation and optimization errors are separately small; it establishes the required full-policy bound for the unchanged policy.

The diagnostic answers a specific method question. The original capped service was less reliable at its tight target, but the saved neural objects were not necessarily inadequate policies. All seven selected policies admit the tighter certificates when verification is extended under a fixed rule. This finding does not alter the original attainment table or justify a comparison between a selectively refined neural certificate and an unrefined comparator as evidence of a work advantage.

\subsection{Direct coupled economic differences}
\begin{table}[htbp]
\centering\small
\caption{Direct neural-minus-ridge cost comparisons, simultaneous 99 percent intervals}
\label{tab:paired}
\begin{tabular}{rrrrr}
\toprule
Price & Comparisons & Positive & Negative & Unresolved\\
\midrule
1 & 12 & 0 & 0 & 12\\
4 & 12 & 0 & 0 & 12\\
\bottomrule
\end{tabular}
\end{table}

For each of the six original neural/ridge pairs, select each method's first $0.25$ crossing. Compare at $(1/8,1/8)$, $(1/4,1/2)$, $(1/2,1/4)$, and $(3/4,3/4)$. Each comparison uses 65,536 common innovation-bin trajectories, 40-bit uniform bins, and the own-policy score in \eqref{eq:score}. The total family error is $0.01$ across 48 one-sided bounds. The complete 24 intervals, endpoint moments, seeds, ambiguity counts, and support bounds appear in the supplement and machine-readable records.

All 24 direct confidence intervals contain zero. Thus the coupled evidence now addresses the actual neural-minus-ridge policy-cost difference, but does not rank the methods. This is stronger information than subtracting regret bounds and different information from a smaller all-state certificate. The simultaneous interpretation pertains to the declared states and fixed policies under the simulation model, not to every state in the compact domain. The all-state policy guarantees remain the separate deterministic statements in Table~\ref{tab:attainment} and Theorem~\ref{thm:policy}.

\subsection{Precision and complete recorded work}
\begin{table}[htbp]
\centering\small
\caption{Original repeated precision catalogue}
\label{tab:precision}
\begin{tabular}{lrrrr}
\toprule
Method & Certified & Accepted32 & Rejected32 & Accepted64\\
\midrule
Fixed32 & 6/36 & 147 & 30 & 0\\
Fixed64 & 36/36 & 0 & 0 & 474\\
Adaptive & 36/36 & 348 & 126 & 126\\
\bottomrule
\end{tabular}
\end{table}

The precision experiment repeats twelve dimension/conditioning/target objects three times. Fixed binary64 and adaptive each certify all 36 services; fixed binary32 certifies six. Fixed64 and adaptive each make 474 accepted updates. Adaptive makes 348 accepted binary32 updates, 126 rejected binary32 proposals, and 126 accepted binary64 updates. Every rejection is charged.

Adaptive is slower in 26 of the 36 matched recorded services and faster in ten; aggregate recorded time is about $0.98\%$ larger. The controller's activation therefore demonstrates its numerical mechanism, not a complete-work saving in this implementation. This result does not invalidate the conditional numerical acceptance theorem. It prevents the theorem from being mistaken for an empirical efficiency conclusion.

\subsection{What the evidence identifies}
The evidence identifies a finite from-primitives nonlinear construction, all-state certificates for its returned policies, a useful retained-policy verification refinement, and direct economic comparisons under the declared state and innovation design. It also identifies the continued strength of conventional comparators. It does not identify a general neural speed advantage, a nonlinear high-dimensional frontier, or a population probability of optimizer success.

An adaptive sparse-grid or adaptive-partition comparator has not been executed for the original neural optimizers. Neither has a nonlinear dimension frontier or repeated isolated timing of every original scalar and coupled service. Section~\ref{sec:catalogue45} adds a separate predeclared horizon/price catalogue with repeated isolated timing for the deterministic constructive backend and spline; it does not replace the original experiments. The retained dimension-explicit analysis exposes the tensor factors those studies would have to confront, rather than treating a two-state experiment as their substitute. The general economic framework and conditional constructive results remain in force; the missing experiments concern comparative evidence for broader implementation claims.

\section{A constructive neural service with a primitive error budget}\label{sec:constructive45}
The policy theorem is deliberately independent of the representation of the continuation. A construction theorem must do more: it must produce a representation, feasible actions, and verified residuals from economic primitives. This section supplies such a theorem for a deterministic ReLU backend. It does not assume convergence of the nonlinear optimizers used in the earlier experiments. Its distinguishing property is that the continuation has a prescribed Lipschitz constant even when its Bellman labels are inexact.

\subsection{Primitive conditions and the neural architecture}
Consider the compact Bellman economy of Section~\ref{sec:economy}, with a common action set $A\subset\mathbb R^a$ and state set $X\subset\mathbb R^d$. All distances in this section use the $\ell^1$ norm. The innovation law and the monotone, cash-invariant functional $\rho_t$ do not depend on the state or action. The transition remains in $X$. Suppose known nonnegative constants satisfy, uniformly in the innovation,
\begin{align}\label{eq:primitive45}
 |c_t(x,a)-c_t(y,b)|&\le C_{x,t}\|x-y\|_1+C_{a,t}\|a-b\|_1,\\
 \|F_t(x,a,z)-F_t(y,b,z)\|_1&\le F_{x,t}\|x-y\|_1+F_{a,t}\|a-b\|_1.
\end{align}
The terminal cost is $L_g$-Lipschitz. These assumptions are additional conditions on the constructive backend, not a replacement for the more general feasible correspondences in the original economic formulation. A state-dependent feasible correspondence requires its own feasible-cover and transport constants.

We require an effective numerical description of the primitives: certified finite state and action covers can be obtained at every positive radius; their nodes are feasible; and a numerical evaluator can enclose each Bellman query to any requested positive tolerance in finite work. The last condition includes evaluation of $\rho_t$, not just evaluation of the network. For expectation with a compact innovation domain, known cell probabilities and a certified innovation cover provide one such evaluator. If $F_t$ is $F_{z,t}$-Lipschitz in the innovation, replacing each cell by a representative at distance at most $h_z$ contributes at most $\beta_tL_{t+1}F_{z,t}h_z$. Arithmetic, probability-mass approximation, and any remaining integration error must be added. The experiment below instead integrates a piecewise-linear network analytically under the original continuous uniform law.

Set
\begin{equation}\label{eq:constants45}
 L_T=L_g,\qquad L_t=C_{x,t}+\beta_tF_{x,t}L_{t+1},\qquad
 L^a_t=C_{a,t}+\beta_tF_{a,t}L_{t+1}.
\end{equation}
Choose state nodes $x_{t,i}$ of covering radius $h_t$ and action nodes $a_{t,j}$ of covering radius $k_t$. Given the continuation already constructed at date $t+1$, compute numbers $\widehat q_{t,ij}$ satisfying
\begin{equation}\label{eq:oracle45}
 |\widehat q_{t,ij}-Q_t(f_{t+1};x_{t,i},a_{t,j})|\le e_t.
\end{equation}
Define $y_{t,i}=\min_j\widehat q_{t,ij}$, choose a minimizing action index $j(t,i)$ with a fixed tie rule, and construct
\begin{equation}\label{eq:network45}
 f_t(x)=\min_i\{y_{t,i}+L_t\|x-x_{t,i}\|_1\}.
\end{equation}
At date $T$, use labels within $e_T$ of $g$ and the same formula with constant $L_g$. At any other date, deploy the action $a_{t,j(t,i(x))}$ attached to a nearest state node, again with a fixed tie rule. The action is feasible by construction. No future optimal value or conventional comparator is used as a label.

Formula~\eqref{eq:network45} is an exact ReLU construction, not an appeal to a universal approximation theorem. Absolute values use $|u|=\max\{u,0\}+\max\{-u,0\}$; minima use $\min\{u,v\}=u-\max\{u-v,0\}$. A balanced tree of these gates implements the finite minimum. With $K_t$ state nodes, the resulting sparse affine--ReLU circuit has $O(dK_t)$ coefficients and $O(1+\log K_t)$ depth with unrestricted affine fan-in. Signed intermediate values can be carried by their positive and negative parts. The construction is related to the classical Lipschitz extension formula of \citet{mcshane1934}; neither that formula nor its elementary ReLU realization is claimed as new. Here its role is to close an own-future dynamic construction, including the deployed actor and every numerical allowance.

\begin{theorem}[Constructive own-future neural policy]\label{thm:constructive45}
Under the preceding conditions, the network~\eqref{eq:network45} is $L_t$-Lipschitz at every date, independently of the numerical label errors. Its residual and its deployed actor satisfy
\begin{align}\label{eq:residual45}
 -e_t\ &\le f_t-T_tf_{t+1}\le L^a_tk_t+e_t+2L_th_t,\\
 Q_t(f_{t+1};x,\pi_t(x))-T_tf_{t+1}(x)&\le L^a_tk_t+2e_t+2L_th_t.\label{eq:actor45}
\end{align}
At the terminal date, $-e_T\le f_T-g\le e_T+2L_gh_T$. Consequently the cost of the returned policy, relative to the optimum over all admissible adapted policies, is uniformly bounded by
\begin{equation}\label{eq:gap45}
 \mathcal G_{\rm con}=
 \sum_{t<T}B_t\{2L^a_tk_t+4e_t+4L_th_t\}
       +B_T\{2e_T+2L_gh_T\}.
\end{equation}
For fixed nodes and $L_t$, changing every label by at most $\delta$ changes the represented function by at most $\delta$ in the uniform norm; it does not change the Lipschitz bound.
\end{theorem}

\begin{proof}
Monotonicity and cash invariance imply nonexpansiveness in the uniform norm. Thus $Q_t(f_{t+1};\cdot,a)$ and $T_tf_{t+1}$ are $L_t$-Lipschitz, while the action modulus is $L^a_t$. At a node, minimization of the numerical queries gives
$-e_t\le y_{t,i}-T_tf_{t+1}(x_{t,i})\le L^a_tk_t+e_t$.
Each cone in~\eqref{eq:network45} is $L_t$-Lipschitz, as is their finite minimum. The lower residual bound follows from the triangle inequality for every cone; the upper bound uses a nearest node twice, once for the cone and once for the Bellman value. At that node the selected action is within $L^a_tk_t+2e_t$ of the continuous-action infimum. Transport to $x$ costs at most $2L_th_t$. The terminal argument is identical without action minimization. Substitution of these widths and actor allowances into Theorem~\ref{thm:policy} gives~\eqref{eq:gap45}. Finally, minima of two arrays whose entries differ by at most $\delta$ differ by at most $\delta$. The supplement gives the construction, measurability, and work details.
\end{proof}

\begin{corollary}[Deterministic termination at an economic tolerance]\label{cor:termination45}
Let $\varepsilon>0$, $H=\sum_{t=0}^TB_t$, and $s=\varepsilon/H$. At nonterminal dates choose
\begin{equation}\label{eq:budget45}
 e_t\le s/8,\qquad L^a_tk_t\le s/8,\qquad L_th_t\le s/16.
\end{equation}
At the terminal date choose $e_T\le s/4$ and $L_gh_T\le s/4$. Terms with zero coefficient impose no restriction. Then the finite constructive service returns a feasible neural policy with loss at most $\varepsilon$. No probability law over initializations and no optimizer stopping criterion is needed.
\end{corollary}
The result follows because each discounted bracket in~\eqref{eq:gap45} is at most $s$. Effective compactness and the numerical evaluator give finite work for these choices. Arbitrarily small tolerances may require arbitrary precision; the corollary does not promise termination in fixed binary64. A finite implementation ladder can still return a recorded failure, even though the unbounded constructive algorithm terminates under the stated hypotheses.

The earlier scalar constructive result in the retained paper implemented piecewise-linear interpolation as a ReLU network. Its sufficient allocation controlled a label-error-to-mesh ratio in the propagated slope. Formula~\eqref{eq:network45} works in arbitrary finite state dimension and enforces the slope bound directly. In particular, label and action allowances can be allocated at order $\varepsilon$ without a separate slope-preservation requirement that those allowances be quadratic in the state spacing. This is an improvement of that sufficient proof budget, not a lower bound for all spline methods or evidence that conventional Lipschitz projection is unavailable.

\subsection{What the construction costs}
The construction does not eliminate state coverage. With $K_t$ state nodes, $J_t$ action nodes, $M_t$ innovation representatives, and network evaluation cost $P_{t+1}$, direct target formation requires
\begin{equation}\label{eq:work45}
 \sum_{t<T}K_tJ_t\{M_t(C_{F,t}+P_{t+1})+C_{\rho,t}+C_{c,t}\}
       +O\!\left(\sum_{t=0}^T dK_t\right)
\end{equation}
arithmetic operations under the indicated evaluation model. $C_{\rho,t}$ includes the finite-risk functional and probability processing. Certified integration can replace the innovation sum but its own cost must be charged. Bit lengths, adaptive arithmetic, and nearest-node policy queries are additional explicit costs, not constants suppressed by a dimension claim. A rectangular state cover with radius $h_t$ can require order $h_t^{-d}$ nodes; a rectangular action cover can require order $k_t^{-a}$ actions. The theorem is therefore a dimension-explicit construction, not an assertion that a nonlinear high-dimensional frontier has been executed or that the curse of dimensionality has been removed.

In one state dimension, the finite envelope can be compiled exactly into a piecewise-linear representation. Two passes compute the greatest $L$-Lipschitz minorant of the node labels. Between neighboring nodes the envelope is the minimum of the two adjacent affine cones. At most one additional intersection is needed per interval. Rational knots and values permit exact antiderivatives; an interval evaluator encloses their numerical evaluation. This compilation evaluates the defining network exactly and does not substitute labels from a separately solved Bellman problem. The spline comparator receives the same exact-integration acceleration, so it is not charged for an artificial dense neural evaluation loop.

\subsection{A predeclared matched service catalogue}\label{sec:catalogue45}
The new experiment tests complete services, not only previously failed neural candidates. It uses a further nonlinear investment laboratory with $x\in[0,1]$, $a\in[0,1/4]$, $\beta=15/16$, and independent $z$ uniform on $[-1/32,1/32]$:
\begin{align}\label{eq:economy45}
 F(x,a,z)&=\tfrac1{32}+\tfrac{11}{16}x+\tfrac1{16}x(1-x)+a+z,\\
 c_p(x,a)&=(x-\tfrac{11}{16})^2+pa^2+4a^4+2(\tfrac38-x)_+^2,\\
 g(x)&=2(x-\tfrac{11}{16})^2+2(\tfrac38-x)_+^2.
\end{align}
The invariant domain follows from monotonicity in $x,a,z$ and the two extreme transitions $0$ and $1$. Valid primitive state moduli are $C_x=23/8$, $F_x=3/4$, and $L_g=17/4$; action moduli are $C_a=p/2+1/4$ and $F_a=1$. The downside term penalizes a low capital stock, and the quartic investment cost discourages rapid adjustment. This is a numerical economic laboratory, not an estimated or calibrated quantitative exercise.

The complete design was fixed before execution: horizons $\{2,4\}$, prices $\{1,4\}$, two methods, three isolated repetitions, and the full common state/action ladder $N=A\in\{16,32,64,128,256,512\}$. The targets are $0.5$, $0.25$, and $0.125$ in units of discounted cost. These are policy-loss tolerances, not equivalence margins for a neural-minus-spline contrast. Every rung restarts from terminal primitives and constructs its own future backward. The neural method uses~\eqref{eq:network45}; the spline uses linear interpolation of its own labels and propagates its actual exact slope bound. Both evaluate the same continuous innovation law analytically and round labels to dyadic multiples of $2^{-40}$. All numerical query errors are outwardly enclosed and enter the certificate.

All rungs are retained, including those preceding the first target crossing. The clock to a target includes every preceding rung, target formation, rational compilation, actor construction, verification, and durable checkpoint output. Fresh processes run sequentially on one selected CPU with one numerical-library thread and an identical non-economic warm-up. Imports and warm-up are excluded from the internal clock but included in a separate process clock. CPU frequency is not fixed. Independent policy-cost comparison is not performed in this catalogue and is not assigned a zero cost. Three timing repetitions describe local dispersion, not a population law for neural optimization.

\begin{table}[tbp]
\centering
\caption{First certified state/action resolution in the complete constructive catalogue}\label{tab:constructive45}
\begin{tabular}{lrrrrrr}
\toprule
Generator & $T$ & $p$ & $0.5$ & $0.25$ & $0.125$ & Final gap \\
\midrule
ReLU envelope & 2 & 1 & 64 & 128 & 256 & 0.06132 \\
Spline & 2 & 1 & 64 & 128 & 256 & 0.05619 \\
ReLU envelope & 2 & 4 & 128 & 256 & 512 & 0.06274 \\
Spline & 2 & 4 & 64 & 128 & 256 & 0.05799 \\
ReLU envelope & 4 & 1 & 128 & 256 & 512 & 0.12236 \\
Spline & 4 & 1 & 128 & 256 & 512 & 0.09711 \\
ReLU envelope & 4 & 4 & 256 & 512 & -- & 0.12503 \\
Spline & 4 & 4 & 128 & 256 & 512 & 0.10221 \\
\bottomrule
\end{tabular}
\par\smallskip{\footnotesize Entries are first qualifying $N=A$ on the fixed ladder. Each row represents three identical deterministic policy/certificate arrays, not three independent economic tasks. A dash is a retained failure at the declared cap $N=512$. Final gaps use the cap for both methods.}
\end{table}

\begin{table}[tbp]
\centering
\caption{Complete prefix work at target $0.125$, or at the failed cap}\label{tab:work45}
\begin{tabular}{lrrrrll}
\toprule
Generator & $T,p$ & $N$ & Queries & Median s & Range s & Exit \\
\midrule
ReLU envelope & 2,1 & 256 & 176,586 & 0.3009 & [0.2982, 0.3094] & target \\
Spline & 2,1 & 256 & 176,586 & 0.2144 & [0.2136, 0.2365] & target \\
ReLU envelope & 2,4 & 512 & 702,924 & 0.8450 & [0.8354, 0.8472] & target \\
Spline & 2,4 & 256 & 176,586 & 0.2151 & [0.2119, 0.2153] & target \\
ReLU envelope & 4,1 & 512 & 1,405,848 & 1.5540 & [1.5497, 1.5626] & target \\
Spline & 4,1 & 512 & 1,405,848 & 1.2654 & [1.2341, 1.2745] & target \\
ReLU envelope & 4,4 & 512 & 1,405,848 & 1.5495 & [1.4602, 1.5506] & cap \\
Spline & 4,4 & 512 & 1,405,848 & 1.2705 & [1.2625, 1.2810] & target \\
\bottomrule
\end{tabular}
\par\smallskip{\footnotesize Queries sum all Bellman evaluations on earlier failed rungs and the reported rung. Clocks include construction, verification and durable checkpoints. Ranges contain the three isolated observations. A cap row is not time to successful certification. Independent economic comparison is not included.}
\end{table}

The neural backend certifies all twelve services at $0.5$ and $0.25$, and nine of twelve at $0.125$. The spline certifies all twelve at every target. The three neural failures are repetitions of the same horizon-four, price-four economy: the final gap is approximately $0.12503$, just above $0.125$. They remain failures rather than being rounded into attainment. Across the 33 common successful method/target/repetition comparisons, the spline reaches the target earlier in every recorded observation. Thus the new execution establishes reproducible complete construction and a transparent frontier, not neural cost superiority. The deterministic termination corollary concerns its specified sufficient budget; the predeclared finite ladder need not contain that budget.

The finite catalogue measures the deterministic backend just specified. It neither overwrites the original neural optimizer's attainment array nor turns the seven retained-policy diagnostics into prospective observations. The theorem supplies a method-level termination statement for this backend; the recorded services establish only the achieved targets, costs, and storage on the declared economies. A smaller certificate is not a lower realized policy cost. The new numerical comparison and the earlier direct policy-cost comparisons therefore remain distinct throughout.


\section{Neural continuation with an action witness}\label{sec:witness46}
A constructed continuation contains information about the action that produced each Bellman label. Discarding that information and subsequently controlling an unrelated actor can charge the same approximation twice. We now retain the label's feasible action through the minimum circuit. This changes the returned policy, not merely the description of an old certificate. The construction continues to use the original economic primitives and its own fitted future.

\subsection{A one-sided policy sandwich}
Write $\mathcal T_t^\pi f_{t+1}(x)=Q_t(f_{t+1};x,\pi_t(x))$. Suppose that, uniformly on the invariant domain,
\begin{equation}\label{eq:onesided46}
 f_t-\mathcal T_tf_{t+1}\le u_t,\qquad
 \mathcal T_t^\pi f_{t+1}-f_t\le d_t,
 \qquad \ell_T\le f_T-g\le u_T.
\end{equation}
The quantities $u_t$ and $d_t$ are signed upper bounds. Their separate absolute values need not enter a policy-loss account.

\begin{proposition}[One-sided policy sandwich]\label{prop:onesided46}
Under the monotonicity, cash-invariance and Bellman comparison conditions of Section~\ref{sec:economy}, inequalities~\eqref{eq:onesided46} imply
\begin{equation}\label{eq:jointgap46}
 0\le J_0^\pi(x)-V_0^*(x)
 \le \sum_{t<T}B_t(u_t+d_t)+B_T(u_T-\ell_T).
\end{equation}
The bound is invariant to additive date-specific shifts of the continuation. It holds for the stated recursive evaluations as well as expectation; no simulation identity is required.
\end{proposition}
\begin{proof}
Backward order comparison gives $V_0^*\ge f_0-\sum_{t<T}B_tu_t-B_Tu_T$ and $J_0^\pi\le f_0+\sum_{t<T}B_td_t-B_T\ell_T$. Subtraction proves the result. A shift $f_t\mapsto f_t+b_t$ adds $b_t-\beta_tb_{t+1}$ to $u_t$ and subtracts it from $d_t$, leaving their sum unchanged.
\end{proof}
The centered theorem remains valid: its choice $d_t=\eta_t-\ell_t$ recovers~\eqref{eq:gap}. The new proposition avoids that indirect bound when the selected-policy residual can be certified jointly with the continuation. It is an order comparison, not a new claim of neural superiority.

\subsection{Retaining the minimizing cone's action}
Use the primitive conditions, grids, query tolerances and labels of Section~\ref{sec:constructive45}. Let $a_{t,i}$ denote the action whose numerical query attains $y_{t,i}$. In place of the nearest-node actor, return
\begin{equation}\label{eq:witnessactor46}
 i_t(x)\in\arg\min_i\{y_{t,i}+L_t\|x-x_{t,i}\|_1\},
 \qquad \pi_t^{\rm w}(x)=a_{t,i_t(x)}.
\end{equation}
A fixed tie convention makes this finite selector measurable. The common action set makes each selected action feasible, even when the attaining cone is not centered at a nearest node. The critic remains the exact affine--ReLU circuit. The actor additionally uses the circuit's comparison index; a discontinuous selector is not asserted to be a continuous ReLU output.

\begin{theorem}[Witness-preserving neural construction]\label{thm:witness46}
Under the conditions of Theorem~\ref{thm:constructive45}, the policy~\eqref{eq:witnessactor46} satisfies
\begin{equation}\label{eq:witnessres46}
 Q_t(f_{t+1};x,\pi_t^{\rm w}(x))-f_t(x)\le e_t.
\end{equation}
Consequently its all-state loss relative to the optimum over all admissible adapted policies is at most
\begin{equation}\label{eq:witnessgap46}
 \mathcal G_{\rm w}=
 \sum_{t<T}B_t\{L_t^ak_t+2e_t+2L_th_t\}
       +B_T\{2e_T+2L_gh_T\}.
\end{equation}
The nonterminal allowance is half the separated sufficient allowance in~\eqref{eq:gap45}; the terminal account is unchanged. This comparison of certified bounds does not assert that the two policies have ordered realized costs.
\end{theorem}
\begin{proof}
For the attaining index $i=i_t(x)$, the numerical label and its action obey $Q_t(f_{t+1};x_{t,i},a_{t,i})\le y_{t,i}+e_t$. Transporting this same action to $x$ costs at most $L_t\|x-x_{t,i}\|_1$. Hence its objective is at most $f_t(x)+e_t$. The unchanged optimal-residual upper bound is $u_t=L_t^ak_t+e_t+2L_th_t$. Apply Proposition~\ref{prop:onesided46} with $d_t=e_t$ and the unchanged terminal enclosure.
\end{proof}
The label, the cone and the selected action are one mathematical object in this argument. Retaining only a Lipschitz-regularized label value without an attaining original action index would not justify~\eqref{eq:witnessres46}. The supplement proves that the scalar compiler preserves precisely this information, including dominated labels, rational intersections and ties.

\begin{corollary}[Resource and execution allowances]\label{cor:witness46}
Let $s=\varepsilon/\sum_{t=0}^TB_t$. At nonterminal dates it suffices to choose
\[
 e_t\le s/4,\qquad L_t^ak_t\le s/4,\qquad L_th_t\le s/8,
\]
and retain the terminal allocation of Corollary~\ref{cor:termination45}. Under the same effective-computability hypotheses the construction terminates with loss at most $\varepsilon$.

If the chosen cone exceeds the minimum by at most $\nu_t$, state acquisition error is at most $r_t$ in the stated norm, and a feasible executed action is within $\alpha_t$ of that cone's action, replace $e_t$ in~\eqref{eq:witnessres46} by
\[
 e_t+\nu_t+2L_tr_t+L_t^a\alpha_t.
\]
Add the discounted execution terms to~\eqref{eq:witnessgap46}.
\end{corollary}
The relaxed allocation is a sufficient-budget improvement relative to the earlier proof, not a complexity lower bound for conventional methods. A tensor cover still incurs its dimension factors. Exact state acquisition is the reference experiment; unchecked lookup using rounded switchpoints does not inherit an exact-index guarantee.

\subsection{A strengthened conventional comparison}
For either the original nearest-node cone actor or the spline actor, let $C_i$ be the node's closed nearest-neighbor cell and let $f_t$ be that method's own continuation. The same selected-label argument gives the computable bound
\begin{equation}\label{eq:nearest46}
 d_t=e_t+D_t,\qquad
 D_t=\max_i\sup_{x\in C_i}
       \{y_{t,i}+L_t\|x-x_{t,i}\|_1-f_t(x)\}.
\end{equation}
On the scalar domain, all terms are piecewise affine. Their maximum is attained among the continuation knots, state nodes and cell boundaries; both adjacent actors are included at a boundary. The exact rational maximum supplies a valid all-state allowance without an additional verification mesh. Both conventional interpolation and the nearest-node neural ablation receive this improved account. Thus the experiment does not contrast the new bound only against a deliberately loose baseline. The old separated bound remains recorded, and the smaller of two bounds for the same returned policy may be used.

\subsection{Matched construction, policy and certificate frontiers}\label{sec:study46}
The following matched catalogue is the frozen R46 experiment. Its statements about direct policy cost refer to the evidence then available; Section~\ref{sec:direct47} supplies the new actual-policy comparison, and Section~\ref{sec:constrained47} executes the state-dependent construction. The design was fixed in the repository before execution. It retains the nonlinear continuous-innovation economy~\eqref{eq:economy45}, horizons $2,4$, prices $1,4$, and all six resolutions $16,32,64,128,256,512$. Three methods and three isolated timing repetitions produce 36 services and 216 rungs. Each rung restarts from primitives. The two cone methods have identical critics and node-action arrays; their deployed selectors differ. The spline constructs its own labels. All methods receive the same exact integration and outward numerical arithmetic.

The declared targets are $1/4$, $1/8$ and $1/16$. Every unsuccessful rung and every capped failure is retained. Internal prefix clocks include target construction, witness or nearest-cell compilation, certificates and durable checkpoints; process clocks additionally include imports and the common non-economic warm-up. Runs use one selected CPU and one numerical-library thread, with method order rotated across repetitions and economic cells. Frequency is not controlled. These are local measurements, not a hardware-independent work theorem. Primitive counts and rational storage are recorded separately from clocks, and are not called complete floating-point or bit-operation counts.

\begin{table}[htbp]\centering
\caption{Final all-state policy-loss bounds at $N=512$}\label{tab:witness46}
\begin{tabular}{rrrrr}\toprule
$T$ & $p$ & Witness neural & Nearest neural & Spline \\ \midrule
2 & 1 & 0.034308 & 0.050092 & 0.039641 \\
2 & 4 & 0.035017 & 0.051637 & 0.040957 \\
4 & 1 & 0.064387 & 0.093593 & 0.066369 \\
4 & 4 & 0.065720 & 0.097039 & 0.069979 \\
\bottomrule\end{tabular}
\par\smallskip\footnotesize Bounds are rounded upward in the records; displayed decimals are summaries, not the target test. All three repetitions have identical checkpoints. Both nearest-node comparators use the strengthened one-sided certificate.
\end{table}

All three methods certify $12/12$ services at $1/4$ and at $1/8$, and $6/12$ at $1/16$. The latter failures are all horizon-four services and are retained. The witness bound is strictly tighter than the nearest-node neural ablation in all 24 distinct economy/resolution comparisons, and tighter than the strengthened spline in 23 of 24. The exception is $T=4$, $p=1$, $N=16$, where the witness and spline bounds are approximately 2.06039 and 2.03515. The declared dyadic targets have the same first-crossing resolution for all three methods in every economic cell. Thus the certificate improvement does not yield an earlier first crossing on this particular ladder. The spline is faster in all 30 common successful target/repetition clock comparisons; the witness construction is faster than the nearest-node neural ablation in all 30. These local clocks and tighter regret bounds establish neither policy-cost superiority nor a general neural speed advantage.

The comparisons isolate the effect of carrying the action witness and of using a joint policy certificate. They do not reclassify earlier capped services or claim that the new selector was previously deployed. In particular, an improvement shared by the strengthened nearest-node ablation cannot be attributed solely to witness selection. A tighter regret certificate does not identify the sign of a direct policy-cost contrast. The existing direct expected-cost intervals and adverse precision findings remain unchanged.

\section{State-dependent constraints and feasible witnesses}\label{sec:feasible46}
Investment opportunities often depend on current capital, collateral or remaining capacity. The common-action construction above does not by itself authorize transporting a node's action to a state where that action is inadmissible. We now incorporate a feasible repair into the same neural Bellman construction. The original controlled economy, rather than a relaxed control problem, remains the comparison object.

Use the invariant state domain, numerical-query conditions and common evaluation functional of Section~\ref{sec:constructive45}. Replace its common action set by nonempty compact sets $A_t(x)$ satisfying
\begin{equation}\label{eq:hausdorff46}
 d_H(A_t(x),A_t(z))\le K_t^A\|x-z\|_1,
\end{equation}
where $d_H$ is Hausdorff distance in the action norm. The primitive inequalities~\eqref{eq:primitive45} are now required for pairs of feasible state--action points; neither costs nor dynamics need to be evaluated at infeasible actions. Assume effective finite $k_t$-nets of $A_t(x_{t,i})$ and a feasible Borel repair $\Gamma_{t,i}(x)\in A_t(x)$ of each stored action $a_{t,i}$, with
\begin{equation}\label{eq:repair46}
 \|\Gamma_{t,i}(x)-a_{t,i}\|_1
 \le K_t^A\|x-x_{t,i}\|_1+\zeta_t.
\end{equation}
The nonnegative $\zeta_t$ is a certified displacement allowance, not permission to violate feasibility. Exact nearest-point selection gives $\zeta_t=0$ when available. The effective repair and feasible-net procedures are additional computational inputs, whose work must be charged.

Starting with $L_T=L_g$, define
\begin{align}\label{eq:graphmod46}
 A_t^Q&=C_{x,t}+\beta_tF_{x,t}L_{t+1},\\
 D_t^Q&=C_{a,t}+\beta_tF_{a,t}L_{t+1},\\
 L_t&=A_t^Q+D_t^QK_t^A.\notag
\end{align}
At each state node, minimize the certified queries over its own feasible action net. Retain the original minimizing action with its numerical label $y_{t,i}$. Form the same cone network $f_t(x)=\min_i\{y_{t,i}+L_t\|x-x_{t,i}\|_1\}$. Let $i_t(x)$ attain this minimum and return $\pi_t(x)=\Gamma_{t,i_t(x)}(x)$.

\begin{theorem}[Feasible-witness transport]\label{thm:feasible46}
Under these conditions and the Bellman comparison hypotheses of Section~\ref{sec:economy}, the constructed continuation is $L_t$-Lipschitz, and the returned policy is feasible and measurable. Uniformly on the original state domain,
\begin{align}\label{eq:graphres46}
 -e_t&\le f_t-\mathcal T_tf_{t+1}
       \le D_t^Qk_t+e_t+2L_th_t,\\
 \mathcal T_t^\pi f_{t+1}-f_t&\le e_t+D_t^Q\zeta_t.\notag
\end{align}
Consequently its loss relative to the optimum over all admissible adapted policies is at most
\begin{equation}\label{eq:graphgap46}
 \sum_{t<T}B_t\{D_t^Qk_t+2e_t+2L_th_t+D_t^Q\zeta_t\}
       +B_T(2e_T+2L_gh_T).
\end{equation}
No optimal continuation is supplied to the label generator. The bound concerns the policy after feasible repair, not the possibly infeasible array of unmodified node actions.
\end{theorem}
\begin{proof}
The primitive inequalities imply a modulus $A_t^Q$ in the state and $D_t^Q$ in the action for feasible pairs. Given an action at $x$, the Hausdorff condition supplies an action at $z$ within $K_t^A\|x-z\|_1$. Comparing the two objectives and reversing the roles of $x,z$ shows that $\mathcal T_tf_{t+1}$ is $L_t$-Lipschitz. The feasible node net therefore gives $-e_t\le y_{t,i}-\mathcal T_tf_{t+1}(x_{t,i})\le D_t^Qk_t+e_t$. The cone argument proves the first enclosure. For the attaining index, transport between the two feasible pairs and~\eqref{eq:repair46} give
\[
 Q_t(f_{t+1};x,\Gamma_{t,i}(x))
 \le y_{t,i}+e_t+L_t\|x-x_{t,i}\|_1+D_t^Q\zeta_t.
\]
The cone equals $f_t(x)$. Apply Proposition~\ref{prop:onesided46} and the unchanged terminal enclosure.
\end{proof}

This result closes a constructive route for a class of endogenous constraints within the original model. It does not assume that all feasible correspondences have a finite Hausdorff modulus. With $s=\varepsilon/\sum_{t=0}^TB_t$, sufficient nonterminal allocations are $e_t\le s/8$, $D_t^Qk_t\le s/4$, $L_th_t\le s/8$ and $D_t^Q\zeta_t\le s/4$. Together with the earlier terminal allocation, they imply finite termination at tolerance $\varepsilon$ when all listed operations are effective. Zero coefficients impose no restriction. When exact repair is available, the larger query allowance of Corollary~\ref{cor:witness46} is again sufficient. Setting $K_t^A=\zeta_t=0$ recovers the common-action witness theorem.

\paragraph{A capacity-constrained investment rule.}
For two nonnegative state coordinates on $[0,1]^2$, let $A(x)=[0,(x_1+x_2)/2]$. This is a state-dependent capital-capacity constraint with $K^A=1/2$ in the $\ell^1$ state norm. The repair $\Gamma_i(x)=\min\{a_i,(x_1+x_2)/2\}$ is feasible, satisfies~\eqref{eq:repair46} with $\zeta=0$, and has an elementary affine--ReLU representation for each fixed witness. Thus changing capacity changes the actor and the constructive state modulus through an explicit term $D^Q/2$. The index selector remains a finite comparison procedure. The supplement provides exact rational checks for this example; these are checks of the theorem and admissibility contract, not an additional multidimensional performance study.

Neither feasible-net construction nor repair is free. With $J_{t,i}$ node actions, target-query work is summed over $\sum_iJ_{t,i}$ rather than a common $K_tJ_t$. Each deployed query also incurs the selected repair's work and storage. Under imperfect state acquisition, repair must be certified feasible at the true state, or throughout the acquisition set for a robust implementation. An error allowance for the objective cannot substitute for that feasibility condition.


\section{From acquired states to feasible policies}\label{sec:acquired47}
A policy is implemented at an acquired state, not at an exact real number available without cost. This distinction matters for a witness selector: a small change in the state can change its index discontinuously. The appropriate object is therefore an objective allowance for the selected witness together with feasibility at the true state. A Lipschitz assumption on the actor would replace, rather than solve, this problem.

At date $t$, an observation $\widehat x$ specifies a nonempty acquisition set $\mathcal D_t(\widehat x)\subset X$ containing the true state, with $\|x-\widehat x\|_1\le r_t$ for every $x\in\mathcal D_t(\widehat x)$. Use the feasible-witness data of Section~\ref{sec:feasible46}. Suppose the numerical selector returns an index $i$ such that
\begin{equation}\label{eq:index47}
 y_{t,i}+L_t\|\widehat x-x_{t,i}\|_1\le f_t(\widehat x)+\nu_t.
\end{equation}
The implemented action $\widehat a_t(\widehat x)$ must belong to $A_t(x)$ for every $x$ in the acquisition set. Assume, on that entire set,
\begin{equation}\label{eq:robustrepair47}
 \|\widehat a_t(\widehat x)-a_{t,i}\|_1
 \le K_t^A\|x-x_{t,i}\|_1+\zeta_t^{\rm acq}.
\end{equation}
The observation and selection rules are measurable and causal. Enlarging the history by these observations leaves the stipulated conditional innovation evaluation unchanged; otherwise the full-information comparator must be defined on the enlarged information structure. The allowances are uniform, not fitted to a favorable sample of observation errors.

\begin{theorem}[Acquisition-aware feasible realization]\label{thm:acquired47}
Under the hypotheses of Theorem~\ref{thm:feasible46}, the implemented action satisfies, uniformly over every permitted observation and true state,
\begin{equation}\label{eq:acquiredres47}
 Q_t(f_{t+1};x,\widehat a_t(\widehat x))-f_t(x)
 \le e_t+\nu_t+2L_tr_t+D_t^Q\zeta_t^{\rm acq}.
\end{equation}
Consequently the original full-information optimum bounds the implemented policy's value below, and its loss is bounded above by
\begin{equation}\label{eq:acquiredgap47}
 \sum_{t<T}B_t\{D_t^Qk_t+2e_t+2L_th_t
       +\nu_t+2L_tr_t+D_t^Q\zeta_t^{\rm acq}\}
       +B_T(2e_T+2L_gh_T).
\end{equation}
The result requires no continuity of the index or actor and no probabilistic independence of measurement error. It applies to the common monotone, cash-invariant recursive evaluation under its original comparison assumptions.
\end{theorem}
\begin{proof}
Transport from the original feasible node action to the implemented feasible action gives
\[
 Q_t(f_{t+1};x,\widehat a_t)
 \le y_{t,i}+e_t+L_t\|x-x_{t,i}\|_1+D_t^Q\zeta_t^{\rm acq}.
\]
The cone at $x$ is at most its value at $\widehat x$ plus $L_tr_t$. Apply~\eqref{eq:index47}, then use $f_t(\widehat x)\le f_t(x)+L_tr_t$. This proves~\eqref{eq:acquiredres47}. Backward order comparison, conditional also on the observation history when necessary, proves~\eqref{eq:acquiredgap47}. Feasibility permits comparison with the optimum over all admissible adapted policies. The proof does not compare distances between the two selected actions at nearby states.
\end{proof}

\begin{corollary}[Capacity, measurement, and action quantization]\label{cor:capacity47}
Let $A_t(x)=[0,\kappa_t(x)]$, where $\kappa_t$ is nonnegative and $K_t^A$-Lipschitz. Suppose a computed lower capacity satisfies
\[
 0\le\underline\kappa_t(\widehat x)\le\kappa_t(x),\qquad
 \kappa_t(x)-\underline\kappa_t(\widehat x)\le 2K_t^Ar_t+\chi_t
 \quad(x\in\mathcal D_t(\widehat x)).
\]
For an action spacing $\Delta_t>0$, implement
\[
 \widehat a_t=\Delta_t\left\lfloor
   \frac{\min\{a_{t,i},\underline\kappa_t(\widehat x)\}}{\Delta_t}
 \right\rfloor.
\]
It is feasible on the whole acquisition set and satisfies~\eqref{eq:robustrepair47} with
$\zeta_t^{\rm acq}=2K_t^Ar_t+\chi_t+\Delta_t$.
Exact, unquantized robust repair uses $\Delta_t=0$ and the minimum itself.
\end{corollary}
\begin{proof}
The downward quantization preserves nonnegativity and the lower capacity. Since $a_{t,i}\le\kappa_t(x_{t,i})$, exact repair at $x$ displaces the node action by at most $K_t^A\|x-x_{t,i}\|_1$. Replacing true capacity by its certified lower bound adds at most $2K_t^Ar_t+\chi_t$; downward quantization adds less than $\Delta_t$. Summing proves the displacement inequality.
\end{proof}

\subsection{Ordinary numerical comparisons}
The execution in Section~\ref{sec:constrained47} uses binary64 scores at dyadic acquired states. Stored labels and the coordinate distances are exactly representable at the declared input resolutions. Put $u=2^{-53}$, $\gamma_2=2u/(1-2u)$, $Y_t=\max_i|y_{t,i}|$, and let $\widetilde L_t$ be the represented slope. In state dimension $d$, each computed cone score differs from its exact score by at most
\begin{equation}\label{eq:leaf47}
 \tau_t=d|\widetilde L_t-L_t|+
        \gamma_2\{Y_t+d|\widetilde L_t|\}.
\end{equation}
The first term charges coefficient conversion; the second charges multiplication and addition. A minimum of the computed scores satisfies~\eqref{eq:index47} with $\nu_t=2\tau_t$, including ties. The supplement states the representability and arithmetic conditions and proves the error account. The general theorem also permits other certified score enclosures.

This implementation does not round a rational switchpoint and assume that its original cell is unchanged. For the scalar comparisons, all cells intersecting a state enclosure are retained. For acquired constrained deployment, a numerical index has the objective allowance~\eqref{eq:leaf47} and its action is repaired against the entire acquisition set. Both routes address discontinuity directly.

\subsection{The identical non-neural envelope}
The envelope construction must be separated from the language used to encode it. The classical Lipschitz extension \citep{mcshane1934}, a native minimum computation, and an exact ReLU realization can represent the same function.

\begin{proposition}[Representation control]\label{prop:representation47}
Fix the node locations, labels, slopes, original action witnesses, and selector tie rule. A conventional min-plus implementation and the affine--ReLU implementation of the same envelope return the same mathematical continuation and policy. They have the same exact policy certificate and direct policy cost. A balanced realization uses $O(dK)$ arithmetic and comparison gates for $K$ cones and depth $O(\log K+\log d)$. Merely changing the representation does not improve the covering or Bellman-query factors.
\end{proposition}
\begin{proof}
For every pair of real numbers, $\min\{a,b\}=a-(a-b)_+$. Applying this identity through a balanced reduction tree proves equality of the continuation. Carrying the same attaining index, with the prescribed tie rule, proves equality of the action witness and repaired action. The certificates and costs concern these mathematical objects. The gate count includes the coordinate absolute values and their sums. The supplement separately bounds floating-point gate errors; exact equality is not asserted for unchecked rounded evaluations.
\end{proof}

We execute this non-neural control on all 24 distinct scalar witness checkpoints. It shares the defining data rather than masquerading as an independently trained competitor. The native selector and its ReLU encoding are thus one construction with two representations. The learned hidden-location methods studied earlier are different candidate generators. Their assessment, and that of the deterministic backend, remains a comparison of actual policies and complete computational work.

\section{An oracle-explicit constrained investment economy}\label{sec:constrained47}
We now execute the feasible-witness construction in a coupled two-state economy. Capital affects both future production conditions and current investment capacity. This makes feasible repair part of the economic policy rather than an algebraic example appended to a common-action calculation.

Let $x\in[0,1]^2$, $\beta=15/16$, and $0\le a\le\kappa(x)=1/8+(x_1+x_2)/16$. A common innovation $z$ is continuously uniform on $[-1/32,1/32]$, with opposite exposures in the two transitions:
\begin{align}\label{eq:economy47}
 F_1(x,a,z)&=1/16+x_1/2+x_2/8+x_1(1-x_2)/16+a/2+z,\\
 F_2(x,a,z)&=1/16+x_2/2+x_1/8+x_2(1-x_1)/16+a/4-z.\notag
\end{align}
Define the shortage and imbalance cost
\[
 b(x)=(x_1-x_2)^2/4+2(1/2-x_1-x_2)_+^2.
\]
Stage cost is $\sum_{j=1}^2(x_j-5/8)^2+b(x)+pa^2+4a^4$; terminal cost doubles the squared-target terms and retains $b$. The predeclared cells have $T\in\{2,3\}$ and $p\in\{1,4\}$. These are normalized theoretical economies, not estimated parameters or calibrated welfare magnitudes.

\begin{proposition}[Finite primitives and continuous-law queries]\label{prop:primitive47}
The state domain is invariant under~\eqref{eq:economy47}. In the $\ell^1$ state norm, valid primitive constants are
\begin{equation}\label{eq:moduli47}
 C_x=13/4,\quad L_g=9/2,\quad F_x=11/16,\quad
 F_a=3/4,\quad C_a=p/2+1/4,\quad K^A=1/16.
\end{equation}
The uniform tensor state grid with $N$ subdivisions has covering radius $h=1/N$. At each node, $N+1$ equally spaced feasible actions have covering radius at most $k=1/(8N)$. The repair is $\min\{a_i,\kappa(x)\}$ and has zero extra displacement allowance. If a continuation is $L$-Lipschitz, $M$ uniform midpoint cells enclose its original continuous-law expectation after adding the deterministic remainder $L/(32M)$ to the numerical integration error.
\end{proposition}
\begin{proof}
The primitive formulas give $F_1\in[1/32,29/32]$ and $F_2\in[1/32,27/32]$ by interval bounds. The sums of absolute derivatives in each state column are at most $11/16$, and the action derivative has $\ell^1$ norm $3/4$. Differentiating the costs on the two shortage regions gives~\eqref{eq:moduli47}. Capacity is $1/16$-Lipschitz and at most $1/4$, proving the grid and repair assertions. Within a uniform shock cell, the mean distance to its midpoint is $1/(64M)$. Since the transition has shock modulus two in $\ell^1$, the expectation error is at most $L/(32M)$. Thus quadrature does not change the innovation law to a discrete one.
\end{proof}

The two generators are the feasible cone witness and conventional fitted-value iteration in the tensor bilinear basis. Each works backward against its own fitted future, uses the same feasible action net, and evaluates the same continuous law with the preceding remainder. No optimal value or other method's continuation is supplied as a label. Rounded labels, their valid query errors, slopes, node actions and all policy bounds are stored.

The conventional actor selects a nearest state node and repairs its action at the current capacity. It receives the one-sided account rather than only the older separated certificate. For the bilinear interpolant, the Bellman-residual enclosure is
\[
 -e_t-L_th\le f_t-\mathcal T_tf_{t+1}
       \le D_t^Qk+e_t+L_th.
\]
The interpolation weights are nonnegative and reproduce coordinates. Their weighted mean $\ell^1$ distance to the surrounding vertices is at most $1/N$, which proves the stated extension from the nodal enclosure.

\begin{lemma}[Complete cell-corner verification]\label{lem:corners47}
For the bilinear actor, the selected-policy allowance is $e_t+D_t$, where
\[
 D_t=\max_i\sup_{x\in C_i}
       \{y_{t,i}+L_t\|x-x_i\|_1-f_t(x)\}.
\]
Here $C_i$ is the closed nearest-node cell and both adjacent choices are admitted on a boundary. It suffices to evaluate the at most nine vertices obtained by splitting each such cell along its node coordinates. Outward endpoint evaluation supplies a valid upper bound on $D_t$ without an additional verification mesh.
\end{lemma}
\begin{proof}
On each resulting subrectangle the distance term is affine and the continuation is bilinear. An affine function minus a bilinear function attains its maximum at a vertex: first maximize over one coordinate, then over the other. The label's feasible action, repaired at the queried state, has objective at most $y_{t,i}+e_t+L_t\|x-x_i\|_1$ by feasible transport. Taking the complete maximum proves the allowance.
\end{proof}

\subsection{A cap obtained from economic accuracy}
In this class the previously abstract covering, integration and repair operations are finite routines with explicit costs. Let $S_T=\sum_{t<T}\beta^t$, $a_p=C_x+K^AC_a$ and $b=\beta(F_x+K^AF_a)=705/1024$. Define
\[
 \overline L=\max\{L_g,a_p/(1-b)\},\qquad
 \overline D=C_a+\beta F_a\overline L.
\]
The primitive recursion gives $L_t\le\overline L$ and $D_t^Q\le\overline D$ at every horizon. Suppose that, apart from the quadrature remainder, each nonterminal numerical query has error at most $\xi$ and terminal labels have error at most $\xi_T$. With $M=N$, Theorem~\ref{thm:feasible46} implies
\begin{align}\label{eq:explicitcap47}
 \mathcal G_N&\le C_{T,p}/N+2S_T\xi+2\beta^T\xi_T,\\
 C_{T,p}&=S_T\{\overline D/8+2\overline L+\beta\overline L/16\}
                     +2\beta^TL_g.\notag
\end{align}
Thus a dyadic $N\ge2C_{T,p}/\varepsilon$ and $\xi,\xi_T\le\varepsilon/[4(S_T+\beta^T)]$ suffice for loss at most $\varepsilon$. Acquisition and action quantization add the explicit terms of~\eqref{eq:acquiredgap47}; these must receive part of the same tolerance. Fixed binary64, a fixed label lattice and a finite three-rung study are not asserted to achieve every positive $\varepsilon$.

For this two-state, one-action, one-shock implementation, Bellman queries number $T(N+1)^3$. Each has $N$ midpoint evaluations. Direct cone evaluation examines $(N+1)^2$ cones, giving $O(TN^6)$ primitive arithmetic work. Bilinear future evaluation uses a local patch, giving $O(TN^4)$ work for that implementation. Both store $O(TN^2)$ critic labels and actions. The present batched program additionally holds $O(N^3)$ state--action arrays and $O(BN^2)$ temporary cone entries for fixed batch size $B$. Feasible acquisition and repair add their recorded deployment work, rather than being included at zero cost.

With $p_b$-bit arithmetic and multiplication cost $\mathsf M(p_b)$, the corresponding arithmetic upper bounds acquire a factor $\mathsf M(p_b+T+\log N)$, with coefficient conversion, comparison and serialization charged separately. The precision must also meet the query enclosure requirement. These are algorithmic upper bounds, not lower bounds for all conventional methods. Proposition~\ref{prop:representation47} shows that the identical native envelope has the same leading factors as its ReLU realization. No representation label removes the tensor dependence.

\subsection{Execution and acquired deployment}
The fixed ladder is $N=M=4,8,16$, with three isolated processes per method and economic cell. The 24 services store all 72 rungs, including failed attempts. Targets are $4,2,1,1/2$ in the normalized cost units. They test the construction-to-certificate service; a mathematically valid coarse bound is not automatically an economically adequate tolerance. The full attained bounds and work therefore accompany the target counts.

\begin{table}[htbp]\centering
\caption{Coupled constrained economy: full-prefix work at $N=16$}\label{tab:constrained47}
\begin{tabular}{rrlrrrr}\toprule
$T$ & $p$ & Method & Bound & Queries & Seconds & Time range\\\midrule
2 & 1 & Witness & 2.3148 & 11,534 & 7.045 & [7.003,7.070] \\
2 & 1 & Bilinear FVI & 2.7960 & 11,534 & 0.326 & [0.325,0.352] \\
2 & 4 & Witness & 2.3691 & 11,534 & 6.992 & [6.975,7.188] \\
2 & 4 & Bilinear FVI & 2.8288 & 11,534 & 0.324 & [0.323,0.325] \\
3 & 1 & Witness & 3.3212 & 17,301 & 10.519 & [10.480,10.552] \\
3 & 1 & Bilinear FVI & 3.8915 & 17,301 & 0.490 & [0.487,0.500] \\
3 & 4 & Witness & 3.4107 & 17,301 & 10.482 & [10.425,10.538] \\
3 & 4 & Bilinear FVI & 3.9321 & 17,301 & 0.494 & [0.489,0.520] \\
\bottomrule\end{tabular}
\par\smallskip\begin{minipage}{0.98\linewidth}\footnotesize Query counts and clocks include all three rungs, including failed attempts. Seconds are medians over three isolated process observations; the range reports their minimum and maximum. CPU affinity and library thread counts are fixed, but CPU frequency is not controlled.\end{minipage}\end{table}

Both methods attain the target 4 in all twelve services, at the last rung, and neither attains 2, 1 or $1/2$ within this cap. The witness's final bound is smaller in each economic cell, but bilinear fitted-value iteration is faster in all twelve common successful comparisons. These coarse all-state certificates establish an executed feasible construction, not an economically precise solution at the unachieved tolerances. The sufficient-cap formula states what further resources the mathematical guarantee requires without relabeling these failures.


For each of the twelve distinct saved witness policies we also execute the selector on a $33$-by-$33$ measurement grid, with coordinate acquisition radii $0$, $2^{-12}$ and $2^{-8}$. Exact robust repair and downward action spacing $2^{-12}$ give 72 deployment cases. Capacity is bounded on the entire acquisition box; it is not checked only at its midpoint. The binary64 score allowance is uniform by~\eqref{eq:leaf47}; the grid records exercise implementation and repair activity rather than replacing that proof.

\begin{table}[htbp]\centering
\caption{Acquired-state and quantized-action deployment}\label{tab:acquisition47}
\begin{tabular}{rrrrr}\toprule
Coordinate radius & Action spacing & Active repairs & Quantized & Largest added bound\\\midrule
$0$ & $0$ & 4,817 & 0 & 0.000000 \\
$0$ & $1/4096$ & 4,817 & 0 & 0.004619 \\
$1/4096$ & $0$ & 6,609 & 0 & 0.022513 \\
$1/4096$ & $1/4096$ & 6,609 & 6,609 & 0.027132 \\
$1/256$ & $0$ & 6,609 & 0 & 0.360204 \\
$1/256$ & $1/4096$ & 6,609 & 0 & 0.364823 \\
\bottomrule\end{tabular}
\par\smallskip\begin{minipage}{0.98\linewidth}\footnotesize Each row covers all twelve distinct constrained witness policies and 32,670 measured state--date queries. Feasibility holds on the entire acquisition box. The additional loss is a uniform theoretical allowance for the executed arithmetic, acquisition and action rule, not an estimated average policy cost.\end{minipage}\end{table}

Repair is active rather than a vacuous operation, and downward action quantization is charged explicitly. At exact acquired input the small nonzero numerical selection allowance remains in the exact records even when it rounds to zero in this table. At a nonzero radius the acquisition term is positive; the total bound is not reported as the exact-state certificate.


\section{Policy costs and economic method choice}\label{sec:direct47}
We next evaluate the actual witness policy, not its distance bound relative to an unknown optimum. The policies are the frozen objects at every common first crossing in the original witness catalogue. There are ten such economic-cell and target combinations. Three timing repetitions of an identical checkpoint are not three independent policies.

The primary initial-state law is uniform on $[0,1]$. The additional fixed states $1/8$, $1/2$ and $7/8$ test particular economic starting positions. At each of these forty groups, the three pairwise contrasts concern witness, nearest-cone and spline policies. A fixed design specifies $131{,}072$ common paths, forty-bit innovation bins and family error $0.01$ for all 120 contrasts. The economic comparison uses the same continuous innovations as the original economy.

Let $[\underline z_m,\overline z_m]$ enclose $g-f_T^m$ and let the initial continuation difference have range $[a_0,b_0]$ over the chosen initial-state law. With the one-sided residuals of~\eqref{eq:onesided46}, the score difference of two policies has deterministic support
\begin{align}\label{eq:directsupport47}
 a&=a_0+\sum_{t<T}B_t(-u_t^m-d_t^n)
                  +B_T(\underline z_m-\overline z_n),\\
 b&=b_0+\sum_{t<T}B_t(d_t^m+u_t^n)
                  +B_T(\overline z_m-\underline z_n).\notag
\end{align}
For a common random initial state, one bounds the difference $f_0^m-f_0^n$, not the two separate ranges. Its extrema are obtained exactly from the union of the rational piecewise-affine knots. Conditional expectation and then integration over the initial law extend Theorem~\ref{thm:paired} to this estimand.

Every simulated bin generates lower and upper score endpoints. When a state enclosure meets a selector boundary, the calculation encloses all possible actors and subsequent states. It neither drops the path nor substitutes an unverified midpoint actor. The mean and variance computations have rational roundoff accounts. For the declared family, $\log(4\cdot120/0.01)<11$; a finite positive Taylor sum for $\exp(11)$ certifies this inequality. Using 11 in~\eqref{eq:bernstein} is conservative and avoids an unchecked logarithm evaluation. The supplement gives the complete endpoint construction and arithmetic audit.

\begin{theorem}[Simultaneous policy comparison and adoption]\label{thm:adoption47}
Conditional on the frozen policies and the independent uniform-bin model, the endpoint intervals constructed with~\eqref{eq:directsupport47} and~\eqref{eq:bernstein} cover all declared expected policy-cost differences simultaneously with probability at least $0.99$. If an interval $[L_{mn},U_{mn}]$ has $U_{mn}<0$, policy $m$ has smaller expected cost than $n$ on this event. For a previously specified tolerance $\eta$, $U_{mn}\le\eta$ certifies noninferiority within that tolerance, whereas $[L_{mn},U_{mn}]\subset[-\eta,\eta]$ certifies two-sided proximity.

More generally, fix a reference policy $b$ and known implementation charges $k_m$. Suppose simultaneous intervals $[L_m,U_m]$ enclose $J_m-J_b$ for a finite candidate catalogue, including $[0,0]$ for $b$. Select $\widehat m$ to minimize $U_m+k_m$. On the coverage event,
\begin{equation}\label{eq:adoption47}
 J_{\widehat m}+k_{\widehat m}
 \le \min_m\{J_m+k_m+(U_m-L_m)\}.
\end{equation}
Thus an economic choice made after seeing all intervals retains its declared uncertainty account. The catalogue comparison is within one common economic task and initial-state law, not across different prices, horizons or evaluations.
\end{theorem}
\begin{proof}
The score identity identifies the expected difference. Its interval endpoints are bounded independent observations within each group. Apply the one-sided empirical Bernstein inequality of \citet{maurer2009} to each lower and upper endpoint and take a union bound over 240 tails. The numerical enclosures only enlarge the valid intervals. Independence between different contrasts is unnecessary. On the resulting event, $J_{\widehat m}-J_b\le U_{\widehat m}$ and $U_m\le J_m-J_b+(U_m-L_m)$. Minimize the former upper bound plus the known charge to obtain~\eqref{eq:adoption47}.
\end{proof}

The study fixes $\eta=1/1024$ in its normalized cost units before execution. It is a stated decision tolerance, not an estimated welfare threshold. Zero remains the superiority threshold. An interval merely intersecting zero meets neither criterion automatically. The fixed pseudorandom implementation supplies repeatable inputs; it does not prove the independent sampling model. The theorem states that model explicitly. No claim about initialization risk or the distribution of newly trained policies follows from these frozen-policy comparisons.

\begin{table}[htbp]\centering
\caption{Direct expected-cost comparisons under a predeclared decision tolerance}\label{tab:direct47}
\begin{tabular}{lrrrr}\toprule
Contrast & Groups & Signed & Within $1/1024$ & Largest endpoint\\\midrule
Witness $-$ Cone nearest & 40 & 0 & 40 & 0.0001161 \\
Witness $-$ Spline & 40 & 0 & 40 & 0.0001975 \\
Cone nearest $-$ Spline & 40 & 0 & 40 & 0.0002078 \\
\bottomrule\end{tabular}
\par\smallskip\begin{minipage}{0.98\linewidth}\footnotesize Signed counts exclude intervals containing zero. ``Within'' requires the entire interval inside $[-1/1024,1/1024]$. Largest endpoint is the largest absolute endpoint over the 40 groups. All 120 intervals use one $0.01$ family error account, conditional on the frozen policies and independent-bin model. Each group uses 131,072 common paths.\end{minipage}\end{table}

Every direct interval contains zero, so the experiment identifies no signed policy ranking. Every interval is also contained in the predeclared decision band. It therefore certifies two-sided expected-cost proximity for the specified initial laws at that tolerance, rather than inferring equivalence from a failure to reject equality. The uniform initial-state law is a distributional conclusion, not an all-state assertion. The individual initial states are separate declared estimands.


\subsection{Work as a function of the certified bound}
A coarse target list can hide a genuine difference between certificate frontiers. For the original saved rungs $j$, let $G_{m,j}$ be the valid bound and let $C_{m,j}$ include all work through rung $j$. The complete first-crossing curve is
\begin{equation}\label{eq:frontier47}
 W_m(\varepsilon)=C_{m,j_m(\varepsilon)},\qquad
 j_m(\varepsilon)=\min\{j:G_{m,j}\le\varepsilon\},
\end{equation}
with $W_m=+\infty$ when the cap is exhausted. Its breakpoints are the recorded bounds; both methods' breakpoints partition the entire positive tolerance axis. We reconstruct every interval of this partition, not only intervals favorable to the witness. This is an exact descriptive account of frozen evidence, not a retrospective change to the original targets.

\begin{table}[htbp]\centering
\caption{Complete tolerance-axis partition of the original scalar frontiers}\label{tab:frontier47}
\begin{tabular}{lr}\toprule Relation on a partition interval & Number\\\midrule
Both attain; witness uses fewer queries & 19 \\
Both attain; spline uses fewer queries & 1 \\
Both attain; same prefix queries & 20 \\
Only witness attains within the cap & 4 \\
Only spline attains within the cap & 0 \\
Both exhaust the cap & 4 \\
\bottomrule\end{tabular}\end{table}

The complete curves contain 19 nonempty tolerance intervals on which both methods attain and the witness uses fewer Bellman queries, one favoring the spline, and twenty with identical counts. Four further intervals are attainable only by the witness; four below both final bounds are unattained. Interval counts are not probabilities or measures of economic importance. Above the largest breakpoint in each cell both first-rung counts coincide. The original targets remain unchanged and still have identical crossing resolutions. These descriptive intervals identify where a sharper certificate changes a resource decision without manufacturing new prospective targets.


\subsection{Economic interpretation and the NBO program}
The new evidence connects a feasible implementation to two distinct economic statements. First, the acquisition-aware theorem bounds the loss of a policy that can actually be executed under uncertain state measurement and capacity constraints. Second, the direct experiment certifies proximity of the original scalar policy costs within an explicitly chosen decision tolerance, so that a user of that catalogue may consider implementation charges without treating separate regret bounds as values. These statements are stronger than a certificate comparison alone.

Neither statement implies that a neural representation is indispensable. The identical native envelope is a mathematical control, and the conventional fitted-value and spline implementations remain strong alternatives. The constrained experiment demonstrates two-state continuous-uncertainty construction and nontrivial repair, not a high-dimensional scaling result or an estimated economic application. Its tight targets remain unachieved. The direct cost experiment does not cover the new constrained policies or nonlinear recursive utility; its expectation-specific conclusion is not transferred to those settings. The controlled-diffusion, recursive-preference, endogenous-preference, temporal-self and game developments below retain their own conditions and original role in the NBO program. Their preservation is not counted as a new comparative execution.



\section{Economic applications and conclusion}\label{sec:conclusion}
The original applications concern policy-dependent recursive utility, endogenous preference specifications, temporal selves, controlled diffusion problems, and dynamic games. Their complete economic statements and proofs are retained without alteration. The centered and transport bounds explain how an NBO can support economic counterfactuals when changes in future primitives are controlled. The policy theorem explains what must be verified for the actual policy. Neither statement supplies an equilibrium or transversality condition absent from the application.

The nonlinear investment evidence adds an executable link from primitives to learned continuations and directly certified policies. Its retained-policy diagnostic demonstrates that a failed tight certificate can reflect the resolution of verification rather than the quality of the deployed policy. Its direct coupled comparison answers a different question---the relative cost of two policies---without confusing separate regret bounds with a method contrast.

The acquired-state result makes that returned policy implementable under a stated information and arithmetic contract. The constrained economy supplies finite covering, integration and repair routines; the direct scalar experiment supplies simultaneous tolerance-based cost comparisons and a catalogue-adoption rule. Their distinct assumptions and work accounts remain visible.

Neural Bellman Operators thus organize construction, certification, economic reuse, and resource accounting around the policy actually returned. Their scientific assessment requires all of these objects. Flexible conventional methods remain important competitors; numerical activation is not a work saving, and a fixed catalogue is not a general reliability theorem. Theorem~\ref{thm:constructive45} and Corollary~\ref{cor:termination45} now supply a separate deterministic reliability statement for an explicitly constructed neural backend. They do not retroactively prove convergence of the earlier optimizers. Preserving those distinctions makes the original economic program testable while retaining its full theory and applications.

\appendix
\section{Reading and preservation map}\label{app:map}
The present main article is the authoritative R47 exposition, extending the complete R46 manuscript and retaining its earlier results. The active technical supplement contains complete proofs of its policy, retained-policy, paired expectation, and constructive-neural results. The response maps every B1--B9 and M1--M10 comment of the latest R46 report to a mathematical result, an execution, or a comparative requirement not established by the new evidence. All active R45 labels, including the retained R44 labels, are preserved. Those earlier review responses remain preserved. The current response addresses the R46 report; acquired-state feasibility, explicit constrained construction, direct policy-cost proximity and complete tolerance curves are stated and proved in this revision.

The current source tree also materializes the complete previous R41 article, supplement, and economic-applications companion without changing their bytes. They preserve the full controlled-diffusion, recursive-preference, endogenous-preference, temporal-self, game, and historical numerical developments. Their presence is preservation of those results under their original hypotheses, not new empirical validation across all applications. The current argument and new constructive proofs can be read without selecting another revision branch. Original R42 capped-service outcomes, R44 postselected diagnostics, and the predeclared R45 constructive catalogue are separately identified. No earlier review or revision is overwritten.

\bibliographystyle{ecta-fullname}
\bibliography{references}
\end{document}
